% =====================================================================
%  R. Nielsen, HR. PROFESSOR BRØCHNERS PHILOSOPHISKE KRITIK GJENNEMSEET.
%  Kjøbenhavn. Forlagt af den Gyldendalske Boghandel (F. Hegel).
%  Trykt hos J. H. Schultz. 1868.   52 pp. (text pp. 3--52).
%
%  Diplomatic transcription from a Google Books digitisation of the
%  British Museum (now British Library) copy, shelfmark "8467.bbb.41"
%  (written on the title page), `bibliotek/Nielsen, Rasmus/
%  1868-brochners-kritik.pdf`, 60 PDF pages, 354.72 x 580.32 pt, each
%  body page one 1-bit JBIG2 image at 600 ppi,
%  sha256 6597e922a147fcff515f6019ca31ce6fd35f195f092be01397f42f2fa933104d.
%
%  The middle round of a three-part exchange: Brøchner, Problemet om Tro
%  og Viden (1868) -> this reply -> Brøchner, Et Svar til Professor
%  R. Nielsen (1868; texts/brochner/svar-nielsen), whose p. 3 quotes this
%  title with "gjennemseet af R. Nielsen".
%
%  The exact title, from the title page (PDF 5): "Hr. Professor Brøchners"
%  / "Philosophiske Kritik" (capital P on the title page and wrapper; the
%  text itself writes "philosophiske Kritik", p. 3, p. 52) / "gjennemseet"
%  (letterspaced) / "af" / "R. Nielsen." (letterspaced). No apostrophe in
%  "Brøchners" (Google's metadata adds one).
%
%  This header is the edition's apparatus. The transcription harness
%  (pagemap.py, ocr.sh, spacing.py, check.py, splice.py, joints.py) is not
%  tracked in git and is unrunnable without the scan, so everything an
%  outside reader needs in order to check this text against the scan is
%  recorded here.
%
% ---------------------------------------------------------------------
%  1. PAGE MAP  (verified by eye 22 Sep 2026; see pagemap.py)
% ---------------------------------------------------------------------
%
%      printed 3--52   =   PDF page - 4      (PDF 7--56)
%      title leaf pp. [1]-[2] = PDF 5-6
%
%  Uniform. The head of every one of PDF 7--56 was rendered and its folio
%  read on the image: 4--52 in unbroken sequence on PDF 8--56, no leaf
%  missing or doubled, no misprinted folio. p. 3 (PDF 7) is the opening
%  page and carries no folio. The embedded layer's folios agree throughout.
%
%  Leaves (PDF pages), none of them transcribed except the title page:
%      PDF   1--2   Google Books notices (English, Danish)
%      PDF   3      front wrapper: the title repeated inside a ruled frame
%      PDF   4      inside front wrapper: publisher's advertisement of
%                   Nielsen's own writings ("Paa den Gyldendalske Boghandels
%                   Forlag er udkommet følgende Skrifter af Professor
%                   R. Nielsen"), in Fraktur
%      PDF   5      TITLE PAGE = p. [1] (ms. shelfmark, a pencil mark and
%                   underline at the author's name: not transcribed)
%      PDF   6      p. [2], blank but for the British Museum stamp
%      PDF  56      p. 52: text ends, closing rule, BM accession date stamp
%      PDF  57      advertisement leaf (Brandes, Dualismen i vor nyeste
%                   Philosophie; Hegaard; Nielsen/Hegaard, om Tro og Viden;
%                   Martensen, Den danske Folkekirkes Forfatnings-Spørgsmaal)
%      PDF  58      advertisement (Bjørnson, Fiskerjenten; Ibsen, Peer Gynt;
%                   Welhaven, Samlede Skrifter; "Jesu Christi
%                   Lignelser")
%      PDF  59      blank
%      PDF  60      back wrapper: the PDF 58 advertisement in a second setting
%  There is no Forord, no Indhold, no Rettelser.
%
% ---------------------------------------------------------------------
%  2. STRUCTURE  (from the images, not the OCR)
% ---------------------------------------------------------------------
%
%  One continuous polemic, pp. 3--52, with no heads, no numbered sections
%  and no internal rules. p. 3 opens without a title, on a large initial
%  capital "I et for nylig udgivet Skrift: Problemet om Tro og Viden, ..."
%  (set as ordinary text). p. 52 ends "... i Hr. Prof. Bröchners
%  philosophiske Kritik." followed by a short closing rule.
%  "§ 15" (p. 11) is a citation, not a head.
%
% ---------------------------------------------------------------------
%  3. TYPOGRAPHY AND TRANSCRIPTION CONVENTIONS
% ---------------------------------------------------------------------
%
%  Body in ANTIQUA (roman) throughout -- unlike Brøchner's reply, which is
%  Fraktur. aa (not å). Emphasis is by LETTERSPACING -> \emph{} (e.g. p. 3
%  "Problemet om Tro og Viden", p. 25 "sin Form"); the notes use it too.
%  Latin and other foreign words are set in the same roman as the body
%  ("Non datur metaphysica in mathesin via", p. 25; "instar omnium",
%  p. 51) and are NOT marked; \textit{} only where the image shows a true
%  italic, which set-up did not find.
%
%  The name is printed in TWO forms, Brøchner and Bröchner (the latter
%  with a diaeresis ö, some thirty times, the former about ten), each
%  verified on the image at set-up. Both are transcribed as printed.
%
%  Quotation marks „...“, low-high. Long rule (pause, and page ranges such
%  as "S. 178---195") -> ---. The id-est mark ɔ: (p. 22) is typed U+0254.
%  Footnotes marked *), **) per printed page (p. 25 carries two):
%  \thefootnote gives *), **), ***), reset by \setcounter{footnote}{0}%
%  after every page marker. The p. 50 note runs over onto p. 51; it is set
%  whole at its call with a % comment where the printed page turns.
%  Line-break hyphens are resolved (Cau\-%).
%
% ---------------------------------------------------------------------
%  4. WORD-ACCURACY
% ---------------------------------------------------------------------
%
%  Word-accuracy: every batch diffed against the embedded (Google) text
%  layer before splicing (ocrdiff.py), and again against an independent
%  tesseract -l dan witness. 4 batches; embedded: 7 candidates, 0
%  transcription errors corrected, 0 unresolved; tesseract: 119
%  candidates, 0 corrected, 0 unresolved. See TRANSCRIPTION-PLAYBOOK.md
%  §3 step 4.
%  Because two batches were drafted from the embedded layer, the diff
%  against that layer is not independent of them; each batch agent
%  therefore also read every page of its range at the image against the
%  fragment (48 layer faults removed there: dropped dashes, joined words,
%  a scrambled pair of lines on p. 19, lost quotes and commas).
%  Then a full second collation of ALL 50 pages at the image (200-600
%  ppi, no OCR consulted) by two readers who had not transcribed, 22 Sep
%  2026: 7 errors found and corrected (0.14/page): a silently corrected
%  printer's error (p. 12 Opmærhsomhed), a comma for a semicolon (p. 7),
%  two specks read as full stops (pp. 9, 26), an emphasis extent (p. 29),
%  a spurious paragraph break at the p. 28/29 seam, a word division
%  (p. 48 "der i"). Items the collators left doubtful are kept as
%  transcribed with a % comment at the site (pp. 9, 10, 13, 26).
%  The p. 32 range dash was normalised to --- like pp. 45, 47, 51.
%
% ---------------------------------------------------------------------
%  5. PRINTER'S DEFECTS LOG
% ---------------------------------------------------------------------
%
%  All transcribed as printed, with a % comment at the site:
%  p. 6 "man reddedes" in a quotation where Problemet S. 193 has "men" ·
%  p. 12 Opmærhsomhed (Opmærksomhed) · p. 24 "i „i sin Form“" (doubled i)
%  · p. 24 methaphysisk, p. 25 methaphysiske · p. 25 note "eg" (og) ·
%  p. 27 Piligt (Pligt) · p. 32 Forelæs-/læsninger (dittography across the
%  line break) · p. 36 "jer erkjender" (p. 29 quotes it as "jeg") ·
%  p. 38 kvilket (hvilket) · p. 45 imellen (imellem) · p. 46 phychologisk
%  · p. 47 a „ printed as ",." · p. 49 a ;-shaped mark before „Hr. ·
%  p. 50 note mathemathiske.
%  Not defects: Brøchner/Bröchner both occur and are kept as printed;
%  Væsensvillen (in quotations from Brøchner) is a real form, not a defect.
%  Quote balance „“ is 0 (203/203).
%
% =====================================================================
\documentclass[11pt]{book}
\usepackage[utf8]{inputenc}
\usepackage[T1]{fontenc}
\usepackage{libertinus}
\usepackage{libertinust1math}
\usepackage{textalpha}
\usepackage{graphicx}
\DeclareUnicodeCharacter{0254}{\reflectbox{c}}
\usepackage{amsmath}
\usepackage[danish]{babel}
\usepackage{fancyhdr}
\setlength{\headheight}{28pt}
\addtolength{\topmargin}{-13.5pt}
\pagestyle{fancy}
\fancyhf{}
\fancyhead[LE,RO]{\thepage}
\fancyhead[RE]{\textit{Hr. Professor Brøchners philosophiske Kritik}}
\fancyhead[LO]{\textit{gjennemseet af R. Nielsen}}
\renewcommand{\thefootnote}{\ifcase\value{footnote}\or *)\or **)\or ***)\else
  \arabic{footnote})\fi}
\usepackage{setspace}
\setstretch{1.3}
\usepackage{marginnote}
\newcommand{\opage}[1]{\marginnote{\footnotesize\textit{[#1]}}}
\newcommand{\apage}[1]{\marginnote{\footnotesize\textit{[#1]}}}

\begin{document}

\frontmatter
\begin{titlepage}
\centering
\vspace*{3cm}
{\Large Hr. Professor Brøchners}\\[1em]
{\Huge Philosophiske Kritik}\\[2em]
{\large \emph{gjennemseet}}\\[1.5em]
{\small af}\\[1.5em]
{\large R. \emph{Nielsen}.}\\[5em]
% [publisher's device: roundel with monogram, not transcribed]
Kjøbenhavn.\\[0.5em]
Forlagt af den Gyldendalske Boghandel (F. Hegel).\\[0.3em]
{\footnotesize Trykt hos J. H. Schultz.}\\[0.5em]
1868.
\end{titlepage}

\mainmatter

% --- p. 3 ---
\opage{3}\setcounter{footnote}{0}%
% [p. 3 opens without a head; the first letter "I" is a large initial, set here as an ordinary letter]
I et for nylig udgivet Skrift: \emph{Problemet om
Tro og Viden}, har Hr. Professor Brøchner, efter at
have afhandlet og gjendrevet adskillige andre Theorier,
tilsidst af nogle fra mine Leilighedsskrifter og en Andens
Referat laante Udtalelser sammensat sig et Noget, han
kalder min „Theorie“. Ved at underkaste dette Noget
en udførlig philosophisk Kritik har Hr. Prof. med kjen\-%
delig Selvfølelse, vægtige Grunde og mange Beviser for\-%
visset sig om, at min „Theorie“ ikke duer, og anbefalet sin
% [a dot stands in the right margin, outside the measure, after "sin": a speck, not transcribed]
som den eneste rette. Da jeg nu i denne Anledning kom\-%
mer til at gjennemsee Hr. Professorens philosophiske
Kritik, bliver det altsaa min Opgave at udfinde, hvilke
Bestanddele af den mig tillagte Theorie, jeg kan vedkjende
mig, og hvorledes det forholder sig med de mange Bevi\-%
ser og vægtige Grunde.

For at forstaae en Kritiker, maa man fremfor Alt
mærke sig hans Standpunkt. Det Standpunkt, hvorpaa
Hr. Prof. Br. har stillet sig, er naturligviis „det strengt
videnskabelige“. I hvad han kalder Theoriens Conse\-%
qvenser, med Hensyn til det Psychologiske (S. 178---195),
med Hensyn til det Ethiske (S. 195---199) og med Hen\-%
syn til det Religiøse (S. 199---210), har den kritiserende
Forfatter sammenhobet endeel formeentlige og virkelige
% --- p. 4 ---
\opage{4}\setcounter{footnote}{0}%
Modsigelser, der klarligen skulle bevise, at den af mig
opstillede Troestheorie maa forkastes, naar man anlægger
Videns Maalestok. Hvad er nu herved opnaaet? Netop
hvad jeg ønsker. Langt fra at undergrave min Troes\-%
lære, har Kritiken just bidraget Sit til at stadfæste, hvad
jeg ved enhver Leilighed atter og atter indskærper, at
Troeslæren umulig kan bestaae for Videnskabens Dom\-%
stol, at den derfor hverken kan eller vil underkaste sig
Videns Afgjørelse. Forsaavidt er Hr. Prof., ifølge en
videnskabelig Selvskuffelse, kommen til at give en Stad\-%
fæstelse, hvor han meente at præstere en Gjendrivelse.
Men saa meget jeg end paaskjønner den i Kritikens
Gjendrivelse indeholdte Stadfæstelse, saa eftertrykkelig
maa jeg tillige paatale den „Vilkaarlighed“, hvormed den
kritiserende Forfatter har undladt at tage Hensyn til det
Eiendommelige ved den Theorie, han vil gjendrive. I
den af mig antydede\footnote{Min Religionsphilosophie er endnu ikke udkommen.}
Troeslære bliver Troen jo ingen\-%
lunde opfattet som en paa blind Autoritet og dunkel
Følelse støttet Antagelse af disse eller hine overnaturlige
Kjendsgjerninger, men paa inderlig bevidst Tilegnelse,
personlig Selvforstaaelse; Troen er, ifølge min Opfattelse,
ikke tankeløs, men væsentlig selvtænkende, et oprindeligt
Livsprincip og just derfor et eiendommeligt Tankeprincip.

Det er Troen, ikke som umiddelbar Antagelse, men som
Tankeprincip, jeg har tillagt en egen \emph{Kritik}, idet jeg har
skjelnet skarpt imellem \emph{Troeskritik} og \emph{Videnskritik}.
Hvad har nu Hr. Prof. Brøchners Kritik i denne An\-%
ledning bragt for Lyset? To formelle Indsigelser, af
hvilke den første ophæver sig selv, fordi den er over\-%
flødig, og den sidste falder bort, fordi den er ubegrundet.
Den første (S. 190---91) lyder paa, at Videnskritiken
ikke kan vedkjende sig Troeskritiken, og herom er slet
% --- p. 5 ---
\opage{5}\setcounter{footnote}{0}%
ingen Strid; den anden lyder paa, at Troeskritiken strider
imod Troens Væsen, men den er uden Begrundelse.
Til Beviis for, at den er uden Begrundelse, lade vi Hr.
Prof. selv tale (S. 192---93): „Naar f. Ex. Fortællingen
om Paradiset „med alle deri indeholdte Charakteertræk“,
siges at have væsentlig og sand Virkelighed, saa er for
Viden en „indre“ Virkelighed, en „høiere Virkelighed“,
i hvilken alle i hiin Fortælling indeholdte Charakteertræk
ere bevarede, kun en Phantasievirkelighed, og Troen vil
ikke nøies med en sublimeret Virkelighed, der mangler
den Facticitet, paa hvilken Troen som Tro maa lægge
afgjørende Vægt. Troen fordrer, at det Factiske con\-%
gruerer med det Ideelle, med Troesidealet, at Virkelig\-%
heden i denne Forstand bliver en høiere Virkelighed, men
den protesterer paa det Bestemteste mod hiin Fordobling
af Virkeligheden, der berøver Troens Gjenstand det Histori\-%
skes reale Betydning. Vilde man saaledes lade Christi
Fødsel, Død og Opstandelse foregaae, ikke i den ydre
historiske Virkelighed, i hvilken Menneskelivet har sin
Skueplads, men kun i hiin høiere Virkelighed, til hvilken
de, som gaaende ind under Underets Begreb, efter Theorien
maae høre hen, saa maatte Troen heri see en Tilintet\-%
gjørelse af hele Christendommens Grundlag. Og Viden\-%
skaben vilde kun kunne give den Ret deri, ligesom den
paa egne Vegne overhovedet, medens den protesterer imod
at unddrages den sande Idealitetens Virkelighed, ligesaa
bestemt maa protestere mod en Fordobling, der forfalsker
Idealitetens Virkelighed ved at sætte en Bastardvirkelig\-%
hed, født af Phantasien, istedet, og som ovenikjøbet maa
blive vilkaarlig og inconseqvent. Ethvert Under maatte
faae hiin høiere Virkelighed til sit Element, men af de
Undere, der have et Ydre, tør man kun føre dem,
der kunne reduceres til Visioner, til „Udslag“ af den
Troendes indre Forhold til det Guddommelige, tilbage
% --- p. 6 ---
\opage{6}\setcounter{footnote}{0}%
til den, idetmindste tør man kun gjøre det uforbeholdent
med Hensyn til dem. Saaledes kan f. Ex. Miraklet med
den brændende Tornebusk blive en indre Virkelighed,
men nu saadanne Mirakler som f. Ex. Overgangen over
det røde Hav, eller Vandets Fremspringen af Klippen
ved Moses's Slag? Disse tør man ikke berøve den ydre
Realitets Virkelighed, og hvorfor ikke? Mon af en Troes\-%
grund? Nei, af et inconseqvent, fra Troens Synspunkt
vilkaarligt, rationelt Hensyn til de historiske Kjendsgjer\-%
ninger, der jo tale for, at Jøderne i den ydre Virkelig\-%
heds Forstand kom over det røde Hav, og at de tør\-%
stende Vandrere i Ørkenen ikke maatte nøies med et „in\-%
dre“, „subjectivt“ Vand, eller et „tilkommende“ Vand,
% sic: "man reddedes" as printed; Problemet (S. 193), which is being quoted, has "men reddedes"
man reddedes fra at omkomme ved et i den ydre Virke\-%
ligheds Forstand virkeligt Vand, som Troen endnu søger
i en virkeligt existerende Kilde“\footnote{Som Anmærkning tilføies: „Jfr. 4 Moseb. XX, 11: og de drak
deraf, og deres Qvæg.“ % [the closer is on the image a raised mark of the 9-shaped form; typed “ by convention]
}.

Dette Stykke Kritik er mærkeligt i to Henseender.

For det Første kan da heraf sees, at den Theorie,
Hr. Prof. angriber, ikke er min Theorie. Det Troesbe\-%
greb, han lægger til Grund, er i den Grad afvigende fra
mit, at jeg ikke vil spilde et Ord paa at belyse Afvigel\-%
sen. Den Opfattelse af Underet, han tillægger mig, „con\-%
gruerer“ med min som Fiirkanten med Cirkelen; Enhver,
som blot vil kaste et Blik i mine Skrifter, kan overbe\-%
vise sig derom. Jeg har, idet jeg paa det Bestemteste
opfatter Begrebet Under ikke som Vidensbegreb, men
som Troesbegreb, aldeles ingen Trang eller Tilskyndelse
til at faae den objective Virkelighed, hvor den virkelig
udgjør et uundværligt Moment i Troesphænomenet, bort\-%
forklaret. Det er i Sandhed ikke mig, der har forsøgt
paa at lade enten Israeliterne eller deres Qvæg slukke
% --- p. 7 ---
\opage{7}\setcounter{footnote}{0}%
Tørsten i subjectivt Vand: Hr. Prof. Brøchners „subjec\-%
tive Vand“ udspringer alene fra hans egen „Kilde“.

For det Andet lære vi af ovenstaaende Udtalelse,
hvor umyndig, tankeløs og ubehjælpsom Troen dog er,
naar den forlades af den høie Videnskab. Den stakkels
videnskabsløse Tro har da ikke engang saamegen Dømme\-%
kraft, at den i en bibelsk Fortælling kan gjøre Forskjel
imellem ydre og indre Virkelighed, imellem Aand og
Bogstav; hvorfra skulde den vel ogsaa faae Dømmekraft,
naar al Dømmekraft bliver til i Videnskaben og tilhører
Videnskaben? Heraf kan man da nok forklare sig, at
den om Videnskritiken uvidende Tro, hvad Fortællingen
om Paradiset angaaer, maa enten tage Alt bogstaveligt
eller forkaste Alt. Til at skjelne mellem lavere og høiere
Virkelighed, mellem materiel Virkelighed og Aandsvirke\-%
lighed, Virkelighed for Aanden, hører jo sund Forstand;
men hvorfra skulde den til sig selv overladte Tro faae
sund Forstand, naar al Forstand er en Frugt af Viden\-%
skaben og tilhører Videnskaben? Troen som Tro kan
for den objective Videns Øine gjøre; hvad den vil,
den er forloren. Forsøger den ved en Kritik at udskille,
hvad der ikke vedkommer dens Væsen, saa har den jo
allerede forgrebet sig paa „Fornuften“, har med det
Samme begyndt at „rationalisere“ sig selv og forberedt
sin Undergang i Videnskaben. Forsmaaer den al Kritik,
og griber iblinde, hvad der tiltaler dens „dunkle, selviske
Drift“, saa forraader den selv sin medfødte Tankeløshed;
den objective Viden vil, for at sige det kort, Intet høre
om Troestanker: Troen er dum og skal være dum, og
derved bliver det.

Spørge vi, hvad der berettiger den philosophiske
Kritik til saa ubetinget at hævde Videns Eneret i Aan\-%
dens Verden, da maae vi naturligviis, for at finde det
rette Svar, gaae tilbage til „Grundlaget“. Forfatterens
% --- p. 8 ---
\opage{8}\setcounter{footnote}{0}%
eneste Lys er Videnskabens Lys; thi Troen er kun et
Mysterium, og Mysteriet det yderste Mørke. Men naar
Videnskaben uden videre faaer Lov til at afvise enhver
anden Maalestok end sin egen; naar alle Forhand\-%
linger om Tro og Viden fra Begyndelsen af føres ude\-%
lukkende fra Videns Standpunkt, naar Alt afgjøres i Vi\-%
den og ved Viden: hvad Under da, at „Problemet om
Tro og Viden“ falder ud til Videns Fordeel! Vel seer
det ud, som skulde Videnskaben gjennem Kritiken godt\-%
gjøre sin Ret til at opløse Troen, den positive Tro paa
Aabenbaringens Under; men det er et Skin. Videnskaben
beviser ikke sin Ret; den postulerer den kun; thi ethvert
Beviis er et Vidensbeviis, altsaa et Beviis, der forudsæt\-%
ter, hvad der netop allerførst skulde bevises, nemlig, at
Vidensbeviset har Gyldighed paa Troens Omraade:
Grundbeviset mangler, og, da Grundbeviset mangler, be\-%
væger den hele Beviisførelse sig afvexlende mellem petitio
principii og circulus in probando.

Strax i Begyndelsen (S. 130) fremsættes til min
store Overraskelse den Paastand, at jeg, for at hævde
„Modsætningen mellem Existens (Liv) og Viden (Theorie)“
og mellem „Subjectivt og Objectivt“ tager Existensens
Begreb i „en vilkaarlig Dobbeltbetydning“ og derpaa
bygger „en sophistisk Slutning“. --- „\emph{Existensen}“, siges
der, „tages nemlig \emph{først}, for at hævdes som almeen
Forudsætning og nødvendig Betingelse for enhver con\-%
cret Bestemmelse, altsaa ogsaa for Viden, i Betydningen
af \emph{individuel Væren i Almindelighed, og der\-%
efter}, for at der gjennem Begrebet Existens kan hæv\-%
des det Praktiske en Prioritet og en Overvægt, i Betyd\-%
ningen af \emph{Villiesværen} („Existensen er praktisk“).
Uden denne Tilsnigelse vilde Argumentationen, der skal
føre til, at det er muligt at sætte Existensen øverst og dog
beholde Theorien, men umuligt at sætte Theorien øverst
% --- p. 9 ---
\opage{9}\setcounter{footnote}{0}%
og dog beholde Existensen, og videre til at godtgjøre
Troens Superioritet over Viden, slet ikke kunne skride
fremad. Og i Beviset for hiin Sætning indeslutter den
første Præmis: en Bevidsthed kan existere uden at være
theoretisk, men ikke være theoretisk uden at existere,
tillige en Usandhed. Al Væren i Individualitetens Form % [a raised speck, half the size of a full stop and above the base line, stands between "Al" and "Væren"; not transcribed (collation 2026-09-22)]
er Existens; der kan altsaa gives en saadan Existens,
der ikke er Bevidsthed, ikke theoretisk. Men naar
\emph{Bevidsthed} bliver Existensens Subject, saa \emph{indeslutter
Existens Theorien}, thi \emph{Bevidstheden er kun exi\-%
sterende som værende theoretisk}“. Tot verba...,
vi skulle nu see, hvad der ligger i disse verba!

At jeg virkelig har brugt og endnu fremdeles agter
at bruge Ordet Existens ikke blot om den lavere, ube\-%
vidste, men ogsaa om den høiere, selvbevidste Tilværen
i „Betydningen af individuel Væren i Almindelighed“,
indrømmer jeg uden Forbehold; men hvorledes der i en
med almindelig Sprogbrug overeensstemmende Anvendelse
af et almindelig bekjendt Ord skal kunne skjule sig enten
„Tvetydighed“ eller „Uklarhed“ eller „Tilsnigelse“,
formaaer jeg ikke at indsee. Overalt, hvor Modsætningen
mellem \emph{Existens} og \emph{Theorie} omhandles i mine Skrifter,
bliver Ordet \emph{Theorie} udtrykkelig taget eenstydigt med
Ordet \emph{Videnskab, videnskabelig objectiv Viden}.
Naar jeg siger, at en Bevidsthed kan existere uden at
være theoretisk, men ikke være theoretisk uden at existere,
saa hører der i Sandhed en egen Art Skarpsindighed til
for at misforstaae Meningen; thi Meningen er, at en
Bevidsthed kan existere uden at være videnskabelig, men
ikke være videnskabelig uden at existere. Men naar
dette uimodsigelig og uden alle Omsvøb er Meningen,
den klare, iøinefaldende Mening, udkræves der da ikke
logiske Briller af usædvanlig Slibning for i en saa simpel
% --- p. 10 ---
\opage{10}\setcounter{footnote}{0}%
Sætning at opdage „Uklarhed“, „Vilkaarlighed“, „Til\-%
snigelse“ og „Sophistik“?

Men, vil man maaskee indvende, Uklarheden bliver
jo beviist, og det allerede paa de første Sider (S. 131 flgd.)
og det med Hensyn saavel paa Opfattelsen af „det Sub\-%
jective“ som af „det Objective“. Vi ville da, førend vi
gaae videre, standse lidt ved disse Kategorier og betragte
dem nærmere hver for sig.

Der indvendes, at „Subjectiviteten“ og „det Sub\-%
jective“ tages i „\emph{Dobbeltbetydning}, saaledes at det
snart kommer til at udtrykke det tilfældigt Individuelle
og det Egoistiske, der falder udenfor det objectivt Gyldige,
snart Subjectivitetens almene Form“. Det oversees, siger
den kritiserende Forfatter, at Subjectivitetens Form og\-%
saa rummer det objectivt Gyldige, at Subjectet optager
det som Norm i sig „uden at Indholdet forandres, at den
subjective Handling kan have objective Principer, at der
er et objectivt Sandt og Gyldigt i den fornuftbestemte
Subjectivitet. Subjectivitetens, Livets Standpunkt kommer
i hiin“ (min!) „Opfattelse til at falde sammen med Ego\-%
ismens Standpunkt“. I denne Fremstilling har Kritiken
indlagt Noget, som ikke vedkommer mig. I min Op\-%
fattelse af Subjectivitetens Væsen falder Eftertrykket
ingenlunde paa Modsætningen mellem det Individuelle
og det Almene, thi enhver existerende Subjectivitet er
som saadan individuel-almeen, men derimod paa Mod\-%
sætningen mellem lavere og høiere Selvhensyn: det lavere
Selvhensyn er \emph{egoistisk}, det høiere er \emph{religiøst}. Paa
Egoismens Trin kommer Individualiteten i mangfoldig
Conflict med Tilværelsens objective Magter; den af
Troens Princip bevægede Subjectivitet vil derimod, idet
den ikke opgiver, men kraftig hævder sit indivi\-%
duelle Væsen, vide sig forsonet saavel med Naturens
som med Aandens guddommelige Love.

% --- p. 11 ---
\opage{11}\setcounter{footnote}{0}%
Ifølge min Synsmaade er det saa langt fra, at det
objectivt Almene skulde være istand til at fortære det
subjectivt Individuelle, at Selvhensynet, det egne, per\-%
sonlige Selvhensyn, tvertimod bliver desto stærkere,
dybere, inderligere, jo høiere Subjectiviteten staaer i Aan\-%
dens Rige. Skulde det virkelig forholde sig saaledes,
som Hr. Prof. Brøchner paastaaer, at „Subjectet optager
det objectivt Gyldige“ --- vel at mærke det i \emph{Viden
og Videnskab} objectivt Gyldige --- „som Norm i sig,
uden at Indholdet forandres, at den subjective Handling
kan have objective“ --- vel at mærke, Videns og Viden\-%
skabens objective --- „Principer“, at „Subjectivitetens og
Livets Standpunkt“ i modsat Fald maa „falde sammen
med Egoismens Standpunkt“: da, ja da følger uimod\-%
sigelig heraf, at den positive Religion maa afskaffes.
Skulde det religiøse Selvhensyn virkelig være Egoisme,
Troen paa en personlig Gud Egoisme, Troen paa Sjælens
Udødelighed Egoisme, Aabenbaringstroen i det Hele Ego\-%
isme, saa kan det jo rigtignok ikke negtes, at Evangeliet, trods
Selvfornegtelse, Taalmod og Kjærlighed, oprindelig er ---
ved „dunkelt Instinkt“ og „selvisk Drift“ et Foster af skjult
Egoisme, og Helten fra Gethsemane den største Egoist,
der har levet paa Jorden.

Som Beviis paa min uklare og feilagtige Opfattelse
af det Objective anføres Følgende (S. 132---33): „Paa
samme Maade som Subjectivitetens Begreb, faaer ogsaa Ob\-%
\emph{jectivitetens} Begreb en \emph{Dobbeltbetydning}. Hvor % ["Ob-" at the line end is set solid; only "jectivitetens" is letterspaced]
Subjectivitetens Hovedstadier udvikles“ --- her henvises
til min philosophiske Propædeutik § 15 --- „betyder det
Objective snart det objective Ydre, eller det blot abstracte
Almene, snart det det Subjectives Totalitet omfattende
Almene og Almeengyldige. Og idet den abstracte Mod\-%
sætning definitivt fastholdes, bliver det vilkaarligt igno\-%
reret, at paa de forskjellige Stadier Subjectets Bestemthed
% --- p. 12 ---
\opage{12}\setcounter{footnote}{0}%
--- hvad der især fremtræder tydeligt paa det Stadium,
der betegnes som den uendelige Subjectivitets, --- netop
er i Kraft af hvad der er bestemt som Tilværelsens ob\-%
jective Princip, og at den udtrykker dette. Det bliver
overseet, at den subjective Sandhed netop er den til\-%
egnede og i Tilegnelsen med sig væsentlig lige objective
Sandhed“.

At den philosophiske Kritiker her paa en temmelig
ukritisk Maade har sammenblandet to aldeles forskjellige
Arter af det Almeengyldige, det subjective og det ob\-%
jective, kan let paavises. Det er, besynderligt nok,
% sic: "Opmærhsomhed" for "Opmærksomhed", as printed (collation 2026-09-22, 600 ppi)
undgaaet hans Opmærhsomhed, at hvad der paa det af
Propædeutiken anførte Sted siges om „den uendelige
Subjectivitet“, ikke sigter til at bevise dens Sandhed,
men til at bevise dens Usandhed. Der staaer: „I kraftig
Eensidighed vælges mellem tvende hinanden modsatte
Udveie. Deels bliver det Objective under alle Skikkelser
saaledes nedsat til et blot Middel, at den nydende Indi\-%
vidualitet kan gjøre sig selv til Tingenes Herre (Hedo\-%
nisme, Aristipp), deels bliver det Objective som det Al\-%
meenfornuftige saaledes indaandet i Bevidstheden, at alt
Endeligt forvandles til et Ligegyldigt; luttret og hærdet i
uendelig Resignation, vil den handlende Subjectivitet bestem\-%
me sig selv, idet den bestemmes af Skjæbnen (Stoicisme; Zeno,
Epictet)“. At herved sigtes paa en ved det Objective frem\-%
kaldt Selvskuffelse, er klart af Sammenhængen, og bliver
til Overflod udtrykkelig sagt i det Næstfølgende, hvor
Charakteristiken af „\emph{den absolute Subjectivitet}“
begynder med følgende Ord: „Den formelle Selvbefrielse % [the opening mark is faint and small on the image; read as „]
ender i det Negative; Subjectiviteten“ --- den Subjectivitet,
der vil bestemme sig i Kraft af Tilværelsens objective
Princip --- „søger Livet, men finder Døden“. Det er
Mere end en Fordreielse af mine Ord, det er en gjennem\-%
gaaende Misforstaaelse af den hele Sag, hvorom her
% --- p. 13 ---
\opage{13}\setcounter{footnote}{0}%
handles. Det bliver, ifølge Hr. Prof. Brøchners For\-%
sikkring, overseet, at den subjective Sandhed --- paa
Hedonismens og Stoicismens Trin --- „netop er den til\-%
egnede og i Tilegnelsen med sig væsentlig lige objective
Sandhed“. Men nu staaer, som bekjendt, det stoiske
Princip i afgjørende Strid med det hedoniske; begge
Principer kunne da umulig være lige objective, med
mindre Principernes objective Gyldighed netop skal
kjendes paa, at de staae i subjectiv Strid med hinanden.
Vistnok holder Stoikeren sig overbeviist om, at det
stoiske Princip er det eneste sande, almeengyldige, at
ethvert fornuftigt Menneske bør leve som Stoiker; men
Hedonikeren er jo fra sin Side som Hedoniker ligesaa
overbeviist om, at Nydelsesprincipet udtrykker „den til\-%
egnede, i Tilegnelsen med sig lige objective Sandhed“,
at det Menneske, der ikke forstaaer at nyde, heller ikke
forstaaer at leve. Udtrykker Hedonikeren „den i Til\-%
egnelsen med sig lige objective Sandhed“, saa udtrykker
Stoikeren den i Tilegnelsen, med sig „lige objective“ % [a small low mark after "Tilegnelsen", read as a comma (the layer's reading); tail indistinct]
Usandhed, og omvendt. Hvo skal nu dømme de stri\-%
dende Parter imellem, og fra hvilket Standpunkt skal
Dommen fældes? „Fra Theoriens, Videnskabens, den
objective Videns Standpunkt,“ svarer den kritiserende
Forfatter. Men dersom den philosophiske Kritiks philo\-%
sophiske Vildfarelse nogensteds kan falde i Øinene, saa
maa det vel være her. En Stoiker kan, som bekjendt,
sætte sig ligesaa sikkert ind i den hedoniske Theorie
som en Hedoniker; omvendt kan en Hedoniker være
ligesaa theoretisk hjemme i den stoiske Theorie som en
Stoiker, og endda bliver Stoikeren ikke Hedoniker,
Hedonikeren ikke Stoiker. Kan der gives et mere
talende Vidnesbyrd mod Hr. Prof. Brøchners „tilegnede,
i Tilegnelsen med sig lige objective Sandhed?“

% --- p. 14 ---
\opage{14}\setcounter{footnote}{0}%
Hvor forskjellige Theorier af denne Natur ere fra
Videnskabens i Sandhed objective Theorier, kan sees af
Exempler i Mængde: Man tænke paa den pythagoræiske
Læresætning: hvorledes skulde det være muligt at sætte
sig ind i Beviserne for denne Læresætning og desuagtet
hævde sig et mathematisk „Standpunkt“, der stod i
aabenbar Strid med det pythagoræiske? Man tænke paa
det kopernikanske System: hvorledes skulde det være
muligt at sætte sig ind i de Beviser, ved hvilke den
nyere Videnskab har godtgjort dette Systems objective
Gyldighed, og endda hævde sig et astronomisk Stand\-%
punkt, der stod i aabenbar Strid med det kopernikanske?
Man tænke paa Optikens Lære om Farverne, paa
Physiologiens Lære om Blodets Kredsløb, kort, hvilke
Videnstheorier, man vil, om dem alle gjælder det, at
klar Indsigt i den objectivt gyldige Theorie falder sam\-%
men med Theoriens Antagelse. Men klar Indsigt i en
subjectiv gyldig Theorie falder, hvad Hr. Prof. Brøchner
saa ganske har overseet, ingenlunde sammen med The\-%
oriens Antagelse. Og hvorfor ikke? Fordi den sub\-%
jective Theories Gyldighed væsentlig beroer paa subjective
Grunde, og de subjective Grunde paa den individuelle
Subjectivitet.

At Stoikeren ikke kan overbevise Hedonikeren, og
Hedonikeren ikke Stoikeren, hvormeget de end gjensidig
arbeide med Theorierne, er en ligefrem Følge af Stand\-%
punkternes praktiske, existentielle Modsætning. Hvad i
al Verden kan det hjælpe, at Stoikeren theoretisk paa\-%
viser de formeentlig objective Modsigelser, hvori Nydel\-%
sestheorien kan hilde den Nydende, saalænge den hedo\-%
niske Subjectivitet vedbliver at sætte Nydelsen som
Livets Høieste. At Modsigelserne ikke kunne løses ob\-%
jectivt i en objectiv Theorie, lærer ham just, at det ikke
er det Theoretiske, men det Praktiske, hvorpaa Efter\-%
% --- p. 15 ---
\opage{15}\setcounter{footnote}{0}%
trykket falder, at det Praktiske væsentlig er subjectivt,
og den subjective Praxis en Nydelsens Kunst. Vistnok
kan Nydelseskunsten heller ikke løse alle Modsigelser,
men hvad beviser saa det? „At Nydelsesprincipet selv er
en Modsigelse!“ Ingenlunde. Det er jo netop Hedoni\-%
kerens uafviselige Grundsætning, at Livsnydelsen selv er
Modsigelsens Løsning, at Modsigelsens Løsning er Ny\-%
delse. Nei, det beviser kun, at Livet frembyder Mod\-%
sigelser, som Ingen, Stoikeren saalidt som Hedonikeren,
formaaer at løse. Naar Livets Modsigelser blive Stoi\-%
keren for stærke, hvad gjør han saa? Han tager sig selv
af Dage. Men den samme, den selvsamme praktiske
Udvei staaer jo ogsaa Hedonikeren aaben. Theorien
kan som Theorie umulig bevirke en Forandring i det
subjective Standpunkt, efterdi Vægten af de Grunde, der
skulle afgjøre den ene Theories Fortrin for den anden,
netop afhænger af Subjectivitetens praktiske Grundsætning.

Kunde der virkelig gives en objectiv, af al individuel
og subjectiv Grundstemning uafhængig Indsigt i den
menneskelige Subjectivitets oprindelige Selvgyldighed, da
maatte det fremfor Alt være Videnskaben magtpaalig\-%
gende, at bringe Spørgsmaalet om Sjælens Udødelighed
til afgjørende Besvarelse; thi at Spørgsmaalet om Sjælens
Udødelighed er uadskilleligt fra Spørgsmaalet om Sub\-%
jectivitetens individuelt almene Realitet, vil vist Ingen
negte. Men at Videnskaben, naar det kommer til Stykket,
ligesaa lidt kan bevise Umuligheden af Sjælens Udøde\-%
lighed, som den er istand til at bevise enten Virkelig\-%
heden eller Nødvendigheden, er en afgjort Sag. Den,
der paastaaer, at Umuligheden kan afgjøres i Kraft af
„Tilværelsens objective Princip“, er ligesaa vist hildet i
et Selvbedrag, som den, der mener, at hans Overbeviis\-%
ning om Nødvendigheden skyldes hans Viden. Den, der
negter Udødeligheden, negter den aabenbart i Kraft af
% --- p. 16 ---
\opage{16}\setcounter{footnote}{0}%
subjectiv, individuel Grundstemning; hans foregivne Vi\-%
denskabsgrunde ere ikke objectivt gyldige Grunde, men
Usandsynligheder, som den subjectivt bestukne Reflexion
meddeler et Skin af objective Beviser. Det gaaer ikke
med en Bedømmelse af Selvets Realitet som med en
Bedømmelse af Regnbuens eller af Skyens eller af Mine\-%
ralets Realitet; Regnbuens, Skyens, Mineralets Realitet
er af objectiv Natur og kan efter Videns almeengyldige
Love bedømmes objectivt; men Selvets Realitet er af
subjectiv Natur og maa bestemmes subjectivt. Det Ob\-%
jectives Realitet paatrænger sig Tanken med Nødvendig\-%
hed; det Subjectives Realitet forkynder sig kun paa Be\-%
tingelse af Frihed. At ville aflede den Realitet, der bør
tillægges Ens eget individuelle, subjective Væsen, af den
Realitet, der ifølge videnskabelig Overveielse maa til\-%
lægges Subjectiviteten i objectiv Almindelighed, er en
bagvendt og selvmodsigende Fremgangsmaade, da Dom\-%
men om den almindelige Subjectivitet netop afhænger af
Dommen om den egne. Den Gyldighed, et Menneske
føler, at han \emph{maa} tillægge den egne Subjectivitet, beroer
i sidste Instans paa den Gyldighed, han \emph{vil} tillægge den;
i denne Forstand gjælder det, at Realitetsbestemmelsen
væsentlig er en Villiesbestemmelse.

Den Erkjendelse, Villien som Villie nødvendig maa
forudsætte, er ikke, hvad Hr. Prof. Bröchner med for\-%
ceret Overlegenhed vil gjøre gjældende, en theoretisk
Erkjendelse, en objectiv Erkjendelse af „Tilværelsens ob\-%
jective Princip“; nei, den er just lige det Modsatte, den
er netop en praktisk Erkjendelse, en subjectiv Erkjen\-%
delse af den egne Tilværelses subjective Princip. Den
er med andre Ord en existentiel Selverkjendelse; men
en existentiel Selverkjendelse kan kun ved vilkaarlig Af\-%
vigelse fra al Sprogbrug kaldes theoretisk.

Den existentielle Selverkjendelse befinde sig nu
% --- p. 17 ---
\opage{17}\setcounter{footnote}{0}%
iøvrigt paa et lavere eller paa et høiere Trin, saa er dog
Saameget vist, at dens videre Udvikling ingenlunde be\-%
tegner dens Overgang til objectiv Viden; tvertimod! jo
renere, klarere og dybere den praktiske Selverkjendelse
er, desto tydeligere adskiller den sig fra den theoretiske.
Paa sit første naturlige Trin er Selverkjendelsen umid\-%
delbart sandselig og den tilsvarende Villie egoistisk; paa
Frihedens aandelige, altsaa paa Subjectivitetens høieste
Trin er Selverkjendelsen ideal, og den tilsvarende Grund\-%
villie religiøs: Troen som personlig Forvisning er ingen
Theorie, den er en indre, oprindelig Aandens Gjerning,
en Act af Erkjendelsens, Selverkjendelsens, den praktiske
Erkjendelses Eenhed med Grundvillien.

Hvorledes det forholder sig med Egoismen og Er\-%
kjendelsen af det Almene, kan nu let oplyses; thi det
forholder sig ingenlunde saaledes, som Hr. Prof. Bröchner
antager, at Erkjendelsen af det objective Almene skulde
kunne frigjøre Subjectiviteten fra dens Egoisme. For\-%
saavidt det Grundmenneskelige oprindelig er i ethvert
Menneske, maa hver enkelt menneskelig Subjectivitet
unegteligt kunne udtrykke det for den menneskelige
Existens Almeengyldige i Forhold til, som Erkjen\-%
delsen klarer sig. Men at dette Almeengyldige ikke
desto mindre er og bliver subjectivt, dets Erkjendelse en
afgjørende subjectiv, praktisk Erkjendelse, er uimodsige\-%
ligt. Den praktiske Erkjendelse er en Villieserkjendelse,
en Erkjendelse, der, idet den betinger Villien, betinges
af Villien, en personlig Selverkjendelse. Saalænge der
er noget Ureent, Driftbundet og Egoistisk i Villien, saa\-%
længe er der ogsaa noget Ureent, Driftbundet og Egoi\-%
stisk i Erkjendelsen, og omvendt: Bevægelsen gaaer i en
psychologisk Cirkel. Den stoiske Egoisme er en anden
end den hedoniske; men Stoikeren er desuagtet ligesaa
væsentlig en Egoist som Hedonikeren. I Videns objec\-%
% --- p. 18 ---
\opage{18}\setcounter{footnote}{0}%
tive Almene abstraheres der vel fra Egoismen, forsaavidt
der i det Hele abstraheres fra det egne Jeg; men denne
Abstraction er ingen personlig Frigjørelse, ingen virke\-%
lig Forløsning: den virkelige Forløsning er kun i Troen.
Al objectiv Viden er betinget af, at den egne tilværende
Subjectivitets personlige Anliggender bringes i Selvfor\-%
glemmelse; det er derfor umuligt, at nogensomhelst ob\-%
jectiv Sandhed skulde formaae at hæve den vidende Sub\-%
jectivitet til aandelig fri Selvstændighed. Hvor den ob\-%
jective Viden begynder, ender Selvhensynet; hvor Selv\-%
hensynet begynder, ender den objective Viden. Med
det tilbagevendende Selvhensyn have vi igjen det Selviske
og med dette den psychologiske Cirkel. Den objective
Erkjendelse kan være reen, og Selverkjendelsen endda
ureen: Befrielsen er kun i Troen. Eet af To, det er
Troens Dilemma: enten maa den ufrie, selviske Subjec\-%
tivitet evig forblive i egoistisk, driftbunden Selverkjen\-%
delse, egoistisk, driftbunden Villie, eller ogsaa paa Be\-%
tingelse af høiere personlig Bistand lade sig forløse af
den oprindelig rene, fra driftbunden Egoisme frigjorte
og just derfor i Aand og Sandhed frigjørende Sub\-%
jectivitet.

Det er en farlig Sag at confundere det Sub\-%
jective med det Objective? Den Troende seer i Aabenba\-%
ringen det Almeengyldige; den Ikketroende fornegter
Aabenbaringen i Kraft af det Almeengyldige. Hvis de
nu begge, den Ikketroende ligesaavel som den Troende,
i selvisk lidenskabelig Iver gaae hen og forvexle det
subjective Almeengyldige med det objective, hvad have
vi saa? Ja saa have vi Fanatismen. Troesfanatismen er
bekjendt nok; Særkjendet er, at den selvsyge Subjec\-%
tivitet i lidenskabelig Forblindelse vil forvandle den sub\-%
jective Sandhed til en objectiv Magt, en Tilværelseslov,
for hvilken alle andre Subjectiviteter ubetinget skulle
% --- p. 19 ---
\opage{19}\setcounter{footnote}{0}%
bøie sig. Vidensfanatismen er mindre bekjendt; men
derfor ikke mindre charakteristisk. Den har naturligviis
ikke sit Udspring fra Videnskaben, men fra en selvisk
Fordreielse af Videnskaben. Istedenfor at tage Erkjen\-%
delsen af det Personlige for hvad den er, en subjectiv,
praktisk Erkjendelse, gjør den Vidende Selverkjendelsen
uden videre objectiv theoretisk. Det er et stort psycho\-%
logisk Selvbedrag. Den Vidende bilder sig ind, at den
Subjectivitet, hvorfra han gaaer ud, er et objectivt al\-%
meengyldigt Væsen, en ubedragelig Metaphysikbevidsthed,
en sand Normalsubjectivitet, medens den dog i Virkelig\-%
heden kun er en Abstraction af den egne individuelle
Subjectivitets tilfældige Tilstand. Langt fra at befries
fra „Egoismen“, bliver den Vidende ved dette Selv\-%
bedrag tvertimod bedaaret af Egoismen; thi at gjøre den
egne Subjectivitets tilfældige Hang til objectiv Maale\-%
stok for alle andre Subjectiviteter, er Egoisme. Hvad
Normalsubjectiviteten „i Kraft af Tilværelsens objective
Princip“, har Trang til, det skal nu hver menneskelig
Subjectivitet væsentlig have Trang til, nemlig at blive
metaphysisk; hvad Normalsubjectiviteten ikke føler Trang
til, en personlig Gud, en Trøst i Bønnen, et evigt Liv
o. s. v., slige Drømmerier og subjective Kjællingerier
tør nu ingen sand --- „objectiv sand“ --- menne\-%
skelig Subjectivitet for Fremtiden føle Trang til. At
føle Trang til Sligt er jo at afvige fra det Normale, at
afvige fra det Normale er at afvige fra det Almene, at
afvige fra det Almene er at afvige fra Sandheden; Maale\-%
stokken er objectiv. Med sin objective Maalestok behøver
Normalsubjectiviteten nu blot --- i al Upersonlighed ---
at faae et lille Anfald af personlig Sandhedsiver, og den
ulmende Vidensfanatisme vil forraade sig ved Angreb
paa Troen.

Af det her Fremsatte maa vel Forskjellen imellem
% --- p. 20 ---
\opage{20}\setcounter{footnote}{0}%
min vitterlige Synsmaade og det Noget, Hr. Prof. Bröchner
har kritiseret under Navn af min „Theorie“, være iøine\-%
faldende. Ved at fastholde denne Forskjel bliver det da
en temmelig let Sag at fælde en begrundet Dom om det
videnskabelige Værd, der i det Hele kan tillægges Hr.
Professorens philosophiske Kritik.

Kritiken har taget sit Udgangspunkt i Psychologien.
Under Overskrift: „\emph{Den psychologiske Modsætning
mellem Villie og Viden}“ meddeles et i taageslørede
Almeensætninger reflecteret Afsnit (S. 133---160), hvis
Øiemed er at vise, hvor skjævt, vilkaarligt og inconse\-%
qvent, jeg under min Fremhæven af den væsentlige For\-%
skjel mellem praktisk og theoretisk Erkjendelse har op\-%
fattet Villiens Forhold til Viden. I den af mig udtalte
Sætning: „det i Subjectiviteten Allersubjectiveste er Selvet;
men det Selviske i Selvet er Selvbestemmen, er Villie,“
skal Grundvildfarelsen stikke. For at rette det Feilgreb
vil Hr. Prof. Bröchner da gaae ud: „\emph{fra Væsensuni\-%
versalitetens} Bestemmelse, der udtrykker sig i den
Forsigværen i den frigjorte Almindeligheds Form, der er
Selvbevidsthedens“, og opfatte „Væsenets Grundbestem\-%
melse som Energie, nærmere som Selvhedsenergie“.
Dersom Dunkelhed og Dybsindighed er Eet og det
Samme, saa er det her anbragte Correctiv unegtelig
meget dybsindigt; dersom taaget Almindelighed er et
Beviis paa objectiv Viden, skal det villig indrømmes, at
dette psychologiske Correctiv er holdt i objectiv Viden.
Men dersom Opgaven er en Sagudvikling, dersom vi,
uden Frygt for at krænke den strenge Videns almene
Fordringer, virkelig tør indlade os paa en Tydeliggjø\-%
relse af eiendommelige Forholds eiendommelige Natur,
da skjønner jeg ikke rettere end, at det med saa stor\-%
artet Eftertryk anbragte Correctiv omtrent falder sammen
med Intet.

% --- p. 21 ---
\opage{21}\setcounter{footnote}{0}%
Hvad her opstilles som Correctiv imod mig, er jo
bogstavelig --- endog indtil Slagordene „Selvhed og Selv\-%
hedsenergie“ --- min egen Anskuelse. Skulde her altsaa
udrettes Noget, da maatte der jo føres et fyldestgjørende
Beviis for, at min Anskuelse var i Strid, ikke med Hr.
Prof. Bröchners eiendommelige Anskuelse, thi han har i
dette Punkt ingen eiendommelig Anskuelse, men i Strid
med sig selv. Er et saadant Beviis nu virkelig ført?
Vi ville undersøge det. Naar jeg siger, at det i Sub\-%
jectiviteten Allersubjectiveste er Selvet, saa har jeg jo
med det Samme angivet en Forskjel imellem Subjectivi\-%
teten og Selvet; naar jeg siger, at det Selviske i Selvet
er Selvbestemmen, er Villie, saa har jeg jo med det
Samme angivet en Forskjel mellem Selvet og det Selviske
i Selvet; naar jeg siger, at det Selviske i Selvet maa,
for at være Villie i egentlig Forstand, være Selvbe\-%
stemmen, saa har jeg jo med det Samme indrømmet en
Forhøielse af Selvbestemmen, en psychologisk Stigen fra
Selvhedsenergiens første dunkle, uvilkaarlige Yttring til
den klare bevidste Concentration, der udtrykkelig mærker
Villien som den sig selv sættende, sig selv hævdende,
paa sig selv beroende Selvhedsact, d. e. som vælgende,
besluttende, fuldbyrdende Villie.

Kunde det Almindelige ved Hjælp af det Alminde\-%
lige i al Almindelighed gjøre Sagen klar, saa maatte
man rigtignok indrømme den kritiserende Forfatter, at
han klarede godt for sig. Thi det skorter saa vist ikke
paa varierende Forsikkringer om Villien i det Almene,
om det Almene i Villien, om Villiens Eenhed med det
Almene i al Almindelighed. „Almindeligheden, der cha\-%
rakteriserer Erkjendelsen“, er „en Almindelighed, der
stræber henimod Concretionen“ o. s. v. (S. 135). „Om\-%
vendt er“ (S. 136) „Enkelthedsbestemmelsen i Villien en
saadan, der med Nødvendighed udfolder det Almene i sig,
% --- p. 22 ---
\opage{22}\setcounter{footnote}{0}%
fordi Villien kun er Villie som det selvbevidste, ɔ: efter
Begrebet almene og universelle, Væsens Villie. Villien
maa udfolde sig fra den ved det Enkelte bestemte, det
Enkelte sættende, Villie til den ved det, Væsenet ud\-%
trykkende, Almeenfornuftige bestemte, det Almene reali\-%
serende Villie.“ Fremdeles: „jo mere Villiesrealisationen
fjerner sig fra det Instinctmæssige, Driftagtige og Egoi\-%
stiske, jo mere bliver den en Virkeliggjørelse af det Al\-%
meenfornuftige, af Væsenets Lov, en Væsensvillen, Ob\-%
jectivitetens, det Almeengyldiges Bestemthed i sig“.
Saaledes bliver det i alle Vendinger, under alle Gjen\-%
tagelser --- ved en petitio principii --- fastsat, at Villiens
Almene er og skal og maa være Videns Almene, „Ob\-%
jectivitetens, det Almeengyldiges Bestemthed i sig“; paa
en saadan petitio principii ere alle videnskabelige For\-%
handlinger spildte.

At en driftig petitio principii maa kunne afføde en
rig Mangfoldighed af tautologiske Cirkelbeviser, er be\-%
gribeligt. Ved Hjælp af det i Almindelighed Almene
bliver det foreløbig beviist, at „den skarpe Modsætning“
mellem theoretisk og praktisk Erkjendelse maa hæves,
for at det, naar denne Modsætning først er hævet, kan
blive rigtig beviist, at det Almene, uden Forskjel mellem
Subjectivt og Objectivt, er det Almene, det Almeenfor\-%
nuftige, det objectivt Almene.

Der er i Hr. Prof. Bröchners Behandling af de Ud\-%
talelser, han kritiserer, en Ubehjælpsomhed, som let kan
bringe en mindre opmærksom Læser bort fra det Hoved\-%
punkt, hvorom Undersøgelsen dreier sig. Om de Kjende\-%
mærker, ved hvilke jeg har søgt at charakterisere Mod\-%
sætningen mellem den praktiske og den theoretiske Er\-%
kjendelse, anføres (S. 137) Følgende. „Den praktiske
Bevidsthed skal være personlig, interesseret, i Conflict
med Existensen, bevæget af Frygt og Haab, fuld af
% --- p. 23 ---
\opage{23}\setcounter{footnote}{0}%
Selvhensyn; den theoretiske upersonlig, contemplativ,
uinteresseret, fortabt i det Almene, bevæget af Nødven\-%
dighed, glemmende sig selv i Begrundelsen af det Ob\-%
jective“. Efter at have opregnet disse Characteermærker
tilføier Hr. Prof. Bröchner saa uden videre: „Den relative
Sandhed i disse Modsætninger er anerkjendt gjennem det
ovenfor Udviklede, men den abstrakte og absolute Mod\-%
sætning er benægtet.“ At den af mig fastsatte Modsæt\-%
ning mellem Villie og Viden er i Kritikens „ovenfor
Udviklede“ benegtet, ligefrem --- ved en Sammenblan\-%
ding af det Almindelige med det Almindelige i al Al\-%
mindelighed --- benegtet, indrømmer jeg; men at der
„gjennem det ovenfor Udviklede“ er fremkommet Noget\-%
somhelst, der kunde antyde endog blot et Forsøg paa
enten at udslette eller dog forviske de anførte Kriterier,
kan jeg umulig indrømme. Det er saa langt fra, at For\-%
fatteren har indladt sig paa nogen særlig Prøvelse af de
særlig angivne Mærker, at han endog lidt længere hen
reent synes at have glemt, at jeg i det Hele har opstillet
saadanne Mærker. Han siger (S. 146): „N. lægger et
Eftertryk paa denne Bestemmelse \emph{praktisk Tænkning},
men den faaer imidlertid hos ham ingensteds en virkelig
Begrebsforklaring, ligesaalidt som det tilsvarende Begreb
\emph{praktisk Erkjendelse}, og det maa overhovedet være
paafaldende, paa hvilken Maade han, under Udviklingen
af sin Theorie, anvender sine psychologiske Kategorier,
uden skarp Begrebsbestemmelse og uden Sondring af
det væsentlig Forskjellige.“

Ved at see, hvordan det hermed hænger sammen,
erholde vi et ikke ubetydeligt Bidrag til videnskabelig
Vurdering af Hr. Professorens philosophiske Kritik.

Som Beviis, og det et afgjørende Beviis, for at der
i Virkeligheden ingen Forskjel er mellem praktisk og
theoretisk Tænkning, anbringes (S. 146) følgende Raison\-%
% --- p. 24 ---
\opage{24}\setcounter{footnote}{0}%
nement. „\emph{Praktisk Tænkning} maatte nærmest betyde
en Reflexion, der bestemmer Sammenhængen mellem
Princip eller Forudsætninger og Conseqventser, mellem
Midler og Øiemed, paa det Praktiskes Gebet, altsaa efter
Forudsætningen, en formel Udviklen af Conseqventser af
det, der ikke opløses i Reflexionen, og denne Reflexion,
der dog overhovedet ikke vilde faae nogen Plads paa
det paradox Ethiskes Gebet, maatte, idet Handlingen
indtræder, tænkes som absorberet, usynliggjort i Hand\-%
lingens Enkelthed. En saadan praktisk Tænken vilde
imidlertid i sin \emph{Form} ikke adskille sig fra den theoretiske
Tænken, der uddrager Conseqventser af et endnu ikke af
Erkjendelsen Assimileret, f. Ex. af et blot empirisk Be\-%
stemt, eller som danner Constructionsbegreber.“ Men
hvad beviser dette? Den, der slutter: Alle Mennesker
ere dødelige, Cajus er et Menneske, altsaa er Cajus
dødelig; han \emph{tænker}; den, der slutter: Alle Mennesker
ere Cylindre, Cajus er et Menneske, altsaa er Cajus en
Cylinder, han \emph{tænker}: den sidste Art af Tænkning er i
% ["er i „i sin Form“": the line-end "i" and the quoted "i" both printed (dittography?); as printed]
„i \emph{sin Form}“ ikke til at adskille fra den første. Men
at Tænkning i \emph{sin Form} altid er Tænkning, var nok
hverken det, der skulde bevises, ei heller det, Forfatteren
vilde bevise.

Skal der ikke lægges Vægt paa Principets og For\-%
udsætningernes Eiendommelighed, da er der ganske vist
ikke Spor af Forskjel mellem theoretisk og praktisk
Tænkning: Tænkning er i Almindelighed altid Tænkning.
Men dersom dette skal være et Beviis mod den væsentlige
Forskjel mellem theoretisk og praktisk Tænkning, da fører
dette os et godt Skridt videre; vi erholde da indenfor
Theorien tillige et Beviis for, at den væsentlige Forskjel
mellem metaphysisk og mathematisk Tænkning er hævet.
% ["methaphysisk" here and "methaphysiske" p. 25 (for metaphysisk): as printed]
Den, der tænker methaphysisk, han tænker; den, der
tænker mathematisk, han tænker: Tænkning er i sin Al\-%
% --- p. 25 ---
\opage{25}\setcounter{footnote}{0}%
mindelighed altid Tænkning; ergo er den mathematiske
Tænkning i \emph{sin Form} ikke til at adskille fra den meta\-%
physiske. Mon Mathematikerne ville hilse denne Op\-%
dagelse som et nyt Fremskridt i Videnskaben?

Den methaphysiske Tænkning er saavel i sin Form
som i sit Indhold ueensartet med den mathematiske.
„Non datur metaphysica in mathesin via“, disse Ord af
en skarpt tænkende Mathematiker\footnote{% ["eg" for og: as printed (image and both witnesses)]
C. Ramus: Differential- eg Integralregning. Fortale.} have et stort Vidnes\-%
byrd i Videnskabernes Historie. „Det er først ved at
fremtræde som mathematisk Videnskab, at Mechaniken,
ligesom de andre Dele af Naturvidenskaben har opnaaet
den Fuldendthed i Formen, den Bestemthed og Evidents,
som er en nødvendig Betingelse for enhver sikker For\-%
udbestemmelse af Phænomenerne som regulerede ved de
givne Kræfter, altsaa for Muligheden af heldige An\-%
vendelser. Man har ofte af Ukyndighed misforstaaet den
mathematiske Analyses Betydning, baade betragtet som
Formen, hvori Undersøgelsen fremtræder, og som Midlet,
hvorved den tilfredsstillende gjennemføres. Det er imid\-%
lertid let at forstaae, at Analysen ikke ved sit eget Ind\-%
hold alene kan bevirke Opdagelsen af Naturkræfternes
Love og Phænomener. Den tjener derimod til at opdage
disse Phænomeners indbyrdes Sammenhæng og deres
Afhængighed af de bestemmende Kræfter\footnote{C. Ramus: Analytisk Mechanik. Fortale.}.“

Den analytiske Mechanik drager virkelig Conseqvenser
af et „blot empirisk Bestemt“, af „et endnu ikke af Er\-%
kjendelsen Assimileret“, forsaavidt den drager Conse\-%
qvenser af visse i Erfaringen givne Kjendsgjerninger; men
denne Omstændighed er saare langt fra, enten at nedsætte
dens objective theoretiske Værdi eller at gjøre dens
% --- p. 26 ---
\opage{26}\setcounter{footnote}{0}%
Tænkning eensartet med „den Reflexion, der bestemmer
Sammenhængen mellem Midler og Øiemed paa det
Praktiskes Gebet“. Hr. Professorens Sammenblanding
af theoretisk og praktisk Reflexion røber overalt en kjen\-%
delig Mangel paa indtrængende Saganalyse.

Et andet, endnu stærkere Tilløb til at sætte over
Kløften mellem den theoretiske og den praktiske Erkjen\-%
delse har Hr. Prof. Brøchner allerede taget S. 139,
hvor han særdeles fremhæver, „at Villiens Enkelthed er
det universelle Væsens Fremtrædelsesform“, og derpaa
vedbliver: „I det saaledes bestemte Enkelte gaaer det
Almene kun paa den Maade til Hvile, at det Enkelte
bliver gjennemsigtigt: Omsætningen er en Omsætning fra
Tænkningens, og udtrykkelig fra den theoretiske Tænk\-%
nings, Almindelighed til den personlige Levens Enkelt\-%
hed; men Individet som fornuftigt villende udtrykker
i Enkeltheden det Almene, giver det Væren i dens Form,
der gjennemsigtig aabenbarer det Almene. Formdiffe\-%
rentsen beherskes altsaa af den væsentlige Eenhed;
Bevidsthedsindholdet bliver under begge Former: som
Erkjendelsesindhold og Villiesindhold, væsentlig det samme.
Jeg erkjender Pligten --- og det er atter: jeg erkjender
den theoretisk --- f. Ex. den Pligt at kæmpe for Fædrelandet;
jeg beslutter mig til at existere i det Erkjendte, og idet Be\-%
slutningen er virkeliggjort, er Indholdet, som realiseret i den
enkelte Villie, væsentlig det samme“. Ihvorvel det hele
Tilløb ligefrem løber ud i en Confusion af det
Almene med det Almene, endog i den Grad, at man i
% ["fristes" and, two lines below, "for.": a small mark after each. Collation 2026-09-22 at 600 ppi: after "fristes" a speck at mid x-height, smaller than a stop -- not transcribed; after "for" a small dot on the base line, doubtful -- kept as printed]
Conseqvens af det opstillede Exempel næsten fristes til
at spørge, hvorledes en Mand, hvem Pligten at kæmpe
for. Fædrelandet er „det universelle Væsens“ Theorie og
„Villiens Enkelthed det universelle Væsens Fremtrædel\-%
sesform“, har kunnet afholde sig fra at gaae med i sidste
Krig: saa skal jeg dog, for at det ikke skal synes, som
% --- p. 27 ---
\opage{27}\setcounter{footnote}{0}%
om jeg tog mig Tingen altfor let, til yderligere Oplys\-%
ning forsøge at vise, hvad Resultatet vilde være blevet,
hvis Hr. Professoren, istedenfor strax at identificere den
praktiske Erkjendelse af Pligten med en theoretisk Er\-%
kjendelse af Videns objective Love, havde taget de af
mig opstillede Forskjelsmærker under nøiere Overveielse
og givet sig Stunder til at anstille en Saganalyse.

For at prøve Forskjelsmærkerne sammenligne man
engang f. Ex. Pligtloven med Tyngdeloven. Ifølge min
Opfattelse er Pligtlovens Erkjendelse væsentlig praktisk,
Erkjendelsen af Tyngdeloven derimod væsentlig theoretisk;
som Love have de jo vistnok begge almeen Gyldighed;
men det Almene i Pligtloven har væsentlig subjectiv, det
Almene i Tyngdeloven væsentlig objectiv Gyldighed.
Hvad dette betyder, kan sees, naar vi bytte Mærkerne
om og gjøre Tyngdeloven subjectiv, Pligtloven objectiv.
Det subjective Moment i Tyngdelovens Erkjendelse ligger
da deri, at den, der skal komme til klar og grundig
Indsigt i Lovens Gyldighed, maa \emph{kunne} tænke og \emph{ville}
tænke, altsaa være en tænkende, villende Subjectivitet;
omvendt fremlyser Pligterkjendelsens objective Moment
deraf, at den, der spørger om Pligten som Pligt i Al\-%
mindelighed, ikke spørger om din eller min Pligt i Sær\-%
deleshed, ikke om, hvad dette eller hiint Subject tilfæl\-%
% ["Piligt" (for Pligt): as printed]
digviis anseer for sin Piligt, men om Pligtens selvgyldige
Væsen. Saaledes have vi da faaet Subjectivitet ind i
Tyngdelovens, Objectivitet ind i Pligtlovens Erkjendelse;
men ere disse Erkjendelser derved blevne eensartede?
Det Hele er tom Formalisme. Saasnart vi betone
det Eiendommelige ved enhver af de tvende Loves Gyl\-%
dighed, maa det Ueensartede falde i Øinene.

% [the tight line prints no visible space in "udelukkende Legemerne"]
Tyngdelovens Gyldighed angaaer udelukkende Legemerne
som Legemer, ikke Sjælene som Sjæle, ikke de tænkende,
villende Subjectiviteter, ikke det i Personlighederne Per\-%
% --- p. 28 ---
\opage{28}\setcounter{footnote}{0}%
sonlige. Loven, der siger os, at alle Legemer ere tunge,
siger jo ikke, enten at alle eller kun nogle Personligheder
ere tunge. Men naar Tyngdeloven slet ikke angaaer
Personlighedens, men alene Legemlighedens Væsen, ikke
har sin Opfyldelse i Personlighedens, men kun i Legem\-%
lighedens \emph{Verden}: saa er den Erkjendende jo herved
udtrykkelig opfordret til at see bort ikke blot fra \emph{sin}
\emph{egen} Personlighed, men fra \emph{al} Personlighed. Med Hen\-%
syn hertil kalder jeg Erkjendelsen af en saa objectiv, uper\-%
sonlig Lov en upersonlig Erkjendelse.

Anderledes med Pligtloven. Erkjendelsen kan be\-%
gynde nok saa abstract, nok saa tilsyneladende objectivt
og theoretisk: det Objective og Theoretiske forvandler
sig øieblikkeligt til et Skin, naar Eftertrykket falder paa
Gyldigheden. For hvilke Existenser har vel Pligtloven
Gyldighed? Ikke for de legemlige, men for de sjælelige,
for de tænkende, villende Subjectiviteter selv, altsaa for
Personligheder. Ingen Enkeltvillie behøver særlig at
indestaae for en samvittighedsfuld Opfyldelse af Tyngde\-%
loven; den skal nok blive opfyldt alligevel. Derimod
har enhver fornuftig Enkeltvillie særlig at indestaae for
en samvittighedsfuld Opfyldelse af Pligtloven; den ene
Enkeltvillie skal ikke i Kraft af nogen objectiv Theorie,
uleilige sig med at indestaae for den anden: Enhver skal
i sin Enkelthed personlig indestaae for sig selv. Som
Personlighedslov er Pligtloven i strengeste Forstand en
Samvittighedslov, Spørgsmaalet om Pligten et Samvittig\-%
hedsspørgsmaal. Er Samvittighedsloven nu en Theorie; er
Samvittighed maaskee en Videnskab? Staaer det fast, at
den existerende Personligheds Viden med sig selv i
Samvittigheden hverken er eller kan være en objectiv
Theorie; staaer det urokkelig fast, at ethvert Pligt\-%
spørgsmaal væsentlig er og maa være et Samvittigheds\-%
spørgsmaal, hvilken Belæring faae vi da ud af Hr. Prof.
% --- p. 29 ---
\opage{29}\setcounter{footnote}{0}%
Bröchners belærende Ord: „Jeg erkjender Pligten, --- og
det er atter: jeg erkjender den theoretisk“?

For at gjøre sig Afstanden mellem det Praktiske og
det Theoretiske i den erkjendende Bevidsthed ret an\-%
skuelig, forsøge man engang at sætte Tyngdeloven paa
lige Trin med Pligtloven og behandle de mechaniske
Problemer som Samvittighedsproblemer; det kan da al\-%
drig feile, at det Urimelige jo maa vise sig. Hvor for\-%
skjellig er ikke den i \emph{det Almene fortabte Viden},
\emph{der bevæges af Nødvendighed}, fra den Erkjendelse,
der betinges af \emph{Selvhensyn} og grunder i \emph{Frihed}!
Da Buffon stred med Clairaut om den Newtonske Lovs
almindelige Gyldighed, faldt det sikkerlig ikke nogen af
de stridende Parter ind at beraabe sig paa et personligt
Selvhensyn og sige til Modstanderen: Jeg føler i min
Samvittighed, at du har Uret. Nei, uden alle Selvhensyn,
uden Hjertesudgydelser eller personlige Forsikkringer,
blev Sagen afgjort objectivt, ved objectivt nødvendige
Grunde. Hvor strengt Videnskaben tager det med Grun\-%
denes objective Nødvendighed, sees blandt Andet deraf,
at Buffon, der holdt paa Newtons Lov og forsaavidt havde
Ret i Sagen, dog i Striden selv fik Uret, idet de Grunde,
hvortil han vilde støtte sig, forkastedes som „blot metaphy\-%
siske og ubrugbare i den exacte Videnskab“. Hvor for\-%
skjellige ere ikke de sande uimodsigelige Vidensgrunde
fra de sande, og dog saa let modsigelige Samvittigheds\-%
grunde! Da Luther stred med Zwingli om Nadver\-%
dogmet, var ingen objectiv, være sig exact eller metaphy\-%
sisk, Theorie istand til at forlige de Stridende, og hvor\-%
for ikke? Fordi Spørgsmaalet oprindelig var, ikke et
theoretisk, men et praktisk Spørgsmaal, ikke et Videns\-%
spørgsmaal, men et Personlighedsspørgsmaal, i strengeste
Forstand et Samvittighedsspørgsmaal: enhver af de stri\-%
dende Parter beraabte sig i sidste Instans, og det med
% --- p. 30 ---
\opage{30}\setcounter{footnote}{0}%
fuld Ret, paa sin Samvittighed. Maaske vil Hr. Prof.
Bröchner her indvende, at de Stridende jo netop opstil\-%
lede hver sin Theorie og just derved selv beviste, at
Samvittighedsspørgsmaalene tilsidst maatte søge deres
høieste Afgjørelse i Theorien; jeg behøver da blot at
minde om, at det paa ingen af Siderne var Theorien,
der bar den personlige Overbeviisning, men omvendt den
personlige Overbeviisning, der bar Theorien. Maaskee
vil Hr. Prof. Bröchner indvende, at slige afgjørende
Selvhensyn høre til et forældet Standpunkt, at Luther
med sin Samvittighed endnu ikke var kommen ud over
„det Dunkle, det Egoistiske og Driftagtige,“ der „hænger
sammen med den individuelle Naturværens uoplukkede
Eenhed,“ til „Videns Almene“, til „Væsensuniversalite\-%
tens Bestemmelse, der udtrykker sig i den Forsigværen
i den frigjorte Almindeligheds Form, der er Selvbevidst\-%
hedens“; jeg behøver da blot at gjøre opmærksom paa,
at en Samvittighedsgrund ligesaa umulig kan forvandles
til en objectiv Vidensgrund, som en Begrundelse af
Tyngdeloven kan forvandles til en Begrundelse af Pligt\-%
loven. Skulde det Vendepunkt i Slægtens Historie no\-%
gensinde indtræde, da „den frigjorte Almindeligheds
Form“ gik i Eet med Theoriens objective Form, saa at
det egne personlige Selvhensyn theoretisk flød ud i
„Væsensuniversalitetens Bestemmelse, der udtrykker sig
i den Forsigværen i den frigjorte Almindeligheds Form,
der er Selvbevidsthedens“, saa maatte det være det
Vendepunkt, da Samvittigheden selv og dermed Pligt\-%
loven blev afskaffet. Maaskee vil Hr. Prof. indvende, at
det „Tilværelsens objective Princip“, der gaaer igjennem
Pligtloven, naturligviis ikke er saaledes objectivt som det
Princip, der gaaer igjennem Tyngdeloven, at der nød\-%
vendig maa gjøres Forskjel imellem Naturnødvendighed
og Frihedsnødvendighed, at han netop selv har indrøm\-%
% --- p. 31 ---
\opage{31}\setcounter{footnote}{0}%
met denne væsentlige Forskjel ved at fremhæve „den
Forsigværen i den frigjorte Almindeligheds Form, der
er Selvbevidsthedens“; jeg behøver da blot at antyde
den Conseqvens, hvormed han i eet Slag slaaer sin egen
Theorie til Jorden. Skal her virkelig lægges Vægt paa
den frigjorte Almindeligheds Form, da bliver jo Forskjellen
mellem den nødvendige og den frigjorte Almindelighed en
gjennemgaaende Principforskjel, altsaa en Sphæreforskjel,
og saa --- er Striden endt.

Naar Hr. Prof. Bröchner vil bekæmpe „den ab\-%
stracte Modsætning mellem det Almene og det Enkelte“,
forsaavidt den skal betinge „Disjunctionen mellem Viden
og Existents“ (S. 132), da har han heri kun Uret, for\-%
saavidt han vil bekæmpe min „Theorie“, men derimod
fuldkommen Ret, forsaavidt han vil bekæmpe sin egen.
Det er, hvad han selv (S. 139) kalder „\emph{Enkeltheds\-%
bestemmelsen} i Villien“, hvorpaa det væsentlig kom\-%
mer an. „Oversees det“, fortsætter Hr. Prof., „ved Op\-%
fattelsen af denne“, nemlig \emph{Enkelthedsbestemmel\-%
sen}, „at Villien kun er Villie som det selvbevidste ɔ:
efter sin Natur universelle og almene Væsens Villie, saa
bliver Villiens Enkelthed og Umiddelbarhedsmomentet i
den ført tilbage til den umiddelbare Naturbestemtheds
Virksomhedsyttring i Driften, den sandselig egoistiske
Attraa, og idet denne nærmest sees fra den Side, gjen\-%
nem hvilken den hænger sammen med den individu\-%
elle Naturværens uoplukkede Enkelthed, der staaer i
umiddelbar Modsætning til Videns Almene, saa bliver den
Virksomhed, der udgaaer derfra, opfattet som overalt bestemt
ved denne Enkelthedens Form, og hvor den ikke overho\-%
vedet unddrager sig Bevidstheden, som den, der idethøieste
gjengives af Bevidstheden som det i sin Umiddelbarhed
med Erkjendelsens Almene Incommensurable. Forlades
derimod denne Opfattelse af Enkeltheden --- der lidet
% --- p. 32 ---
\opage{32}\setcounter{footnote}{0}%
vilde stemme med at Friheden tillægges Villien --- og
fastholdes det, at Villiens Enkelthed er det universelle
Væsens Fremtrædelsesform, er \emph{Aands}bestemthed, bliver
% "Aands" alone letterspaced in "Aandsbestemthed", as printed
Forholdet væsentlig et andet.“

Hvis Theorie, det vel kan være, Hr. Prof. her an\-%
griber, er jeg virkelig ikke istand til at gjætte; kun Saa\-%
meget kan jeg sige, at det er ikke min. Aldrig har det
dog kunnet falde mig ind, at Nogen, der fornuftigviis
maatte antages at kunne læse, vilde være istand til at
bringe en saa mærkværdig Misforstaaelse ud af mine
Skrifter, at jeg virkelig blev anseet for at ville identifi\-%
cere Aandsvilliens \emph{Enkelthed} med „den umiddelbare
Naturbestemtheds Virksomhedsyttring i Driften, den sand\-%
selig egoistiske Attraa“ o. s. fr. Hvormeget en saadan
Opfattelse afviger, ikke blot fra min Synsmaade i det
Hele, men udtrykkelig fra mine allerbestemteste Ud\-%
talelser, hvor langt det er fra mig at ville gjøre Villies\-%
existensen til en dyrisk Enkeltexistens, kan jeg godt\-%
gjøre blot ved at anføre et Par Exempler. „Intet Dyr
kan saaledes fordærve sit Liv og forfeile sit Maal som
Mennesket; thi Dyrets Liv er en Naturgave, Menneskets
en væsentlig Opgave. Som væsentlig Opgave er Livets
Maal et Selvbevidsthedens Formaal, et Frihedens Ideal.
Imod Idealet staaer Virkeligheden; men Virkeligheden,
den stedfindende Naturlighed, maa bearbeides, forandres,
omdannes: her maa ikke blot arbeides, her maa tænkes,
her maa villes, at et Ideal kan virkeliggjøres. Som
den existerende Subjectivitet, der selvvirksom stræber
mod Idealet, er Mennesket væsentlig Aand.“ (Forelæs\-%
% printer's dittography: "Forelæs-" at the line end, "læsninger" on the
% next line, giving "Forelæslæsninger"; transcribed as printed
læsninger ov. phil. Propæd. 1860---61 S. 400). At For\-%
% "1860---61": this dash looks shorter than the pause rule (--- og, top of this page),
% but matches the range dashes on pp. 45, 47, 51, all typed ---; made consistent at collation 2026-09-22
holdet mellem Ideal og Virkelighed oprindelig maa
klares i Erkjendelsen, har jeg da heller ikke overseet;
thi den Fortolkning, Kritiken (S. 144) vil lægge ind i
den Kjendsgjerning, at Villien kan bestemme sig lige\-%
% --- p. 33 ---
\opage{33}\setcounter{footnote}{0}%
gyldig mod eller endog tiltrods for Theorien, „at den
skulde godtgjøre Villiens Utilgjængelighed for Erkjen\-%
delsen“, vedkommer ikke mig. Jeg har saa eftertrykke\-%
ligt som muligt indskærpet Erkjendelsens, forstaaer sig,
den praktiske Erkjendelses, Tænkningens, forstaaer sig,
den praktiske Tænknings Nødvendighed for Villien som
Aandsvillie. „I Forhold til Virkeligheden og det virke\-%
lige Livs Opgaver maa den praktiske Erkjendelse ab\-%
solut have større Oprindelighed end den theoretiske. Vel
maa Villien, for at være Villie, oprindelig forudsætte
Tanken; men deels er det, vel at mærke, ikke den theo\-%
retiske, men den praktiske Tænkning, Villien forudsætter,
deels er Villien med denne sin Forudsætten ligesaa op\-%
rindelig som Tanken. Ligesom der i Tanken er Noget,
der ikke kan afledes af Villien, efterdi Tanken som Tanke
ikke er Villie, saaledes er der ogsaa i Villien Noget, der
ikke kan afledes af Tanken, efterdi Villien som Villie
ikke er Tanke. Vistnok falde saadanne Kategorier som
Valgfrihed, Valg, Beslutning o. fl. ind under Tankens
Belysning, forsaavidt de klare sig i sædelig Erkjendelse,
men i Principet ere de ligefuldt Villieskategorier, deres
Erkjendelse er væsentlig praktisk.“ (Om theor. og prakt.
Erkjendelse S. 7).

Forskjellen mellem den theoretiske og den praktiske
Erkjendelse kan udtrykkelig sees af Forskjellen mellem
Slutning og Beslutning. Den theoretiske Erkjendelse
former sig i \emph{Slutninger}, men ikke i \emph{Beslutnin\-%
ger}: Tanken slutter, men Villien beslutter. Videns ob\-%
jective Slutninger danne sig uden Selvhensyn, i contem\-%
plativ Uinteresserthed, uden Frygt og Haab. Og hvor\-%
for? Fordi de danne sig uden Valgfrihed, uden For\-%
pligtelse, uden personligt Ansvar. I Forhold til Tyngde\-%
loven har den Vidende ingen Valgfrihed, thi de objec\-%
tive Grunde tilstede intet Valg; ingen Forpligtelse, thi
% --- p. 34 ---
\opage{34}\setcounter{footnote}{0}%
Opfyldelsen skeer med Nødvendighed; intet Ansvar, thi
Spørgsmaalet om Brugen angaaer ikke Viden: Viden
hviler i Indsigten. Anderledes med den praktiske Er\-%
kjendelse; dens Problemer ere Villiesproblemer, dens
Grunde Villiesgrunde, dens Slutninger Præmisser for
endelig Beslutning. Den theoretiske Erkjendelse af Plig\-%
ten fører ingenlunde enten til en virkelig Pligtopfyldelse
eller endog til Indsigt i, hvad Pligtopfyldelse vil sige;
den fører tvertimod bort fra begge Dele. „Theorien“
siger, at Enkeltvillien er det Almene, at Pligten lige\-%
ledes er det Almene, og at Villien, idet den beslutter
sig til at opfylde Pligten, viser sig paa een Gang baade
som det Enkelte og det Almene. Det er tomme Ab\-%
stractioner. „Theorien“ kommer, hvormeget der end tales
om det Concrete, aldrig ud over Abstractionen, den Ab\-%
straction, at sætte Villiesbegrebet i Stedet for Villien,
Pligtbegrebet i Stedet for Pligten, Beslutningens Begreb
i Stedet for den virkelige Beslutning. Den virkelige
Beslutning er en indvortes Handling; men just i Hand\-%
lingen concentrerer Villien sig afgjørende til Enkelt\-%
villie med alle en existerende Individualitets Egenheder
og Begrændsninger. Det af Hr. Prof. Bröchner valgte
Exempel, at gaae i Kamp for Fædrelandet, kan oplyse
Sagen. Dersom der for en hvilkensomhelst paagjældende
Enkeltvillie ikke var andet Valg end: enten at gaae i
Kamp eller bryde med Pligtloven, saa kunde der dog til
Nød --- skjøndt Pligtlovens Erkjendelse selv i sin rene
Almindelighed er og bliver subjectiv --- være Tale om
Beslutningens theoretiske Gjennemsigtighed; men at den
Pligt at kæmpe for Fædrelandet er en betinget, en i sin
virkelige Betingethed af særlige Omstændigheder og in\-%
dividuelle Vilkaar afhængig Pligt, sees jo blandt Andet
deraf, at det ikke uden videre er Pligt for Alle, ikke f. Ex.
for Oldinge, Krøblinge, Qvinder og Børn. At den
% --- p. 35 ---
\opage{35}\setcounter{footnote}{0}%
virkelige Betingethed paa ingen mulige Maader lader sig
afgrændse eller bestemme ved theoretisk Erkjendelse,
fremlyser deraf, at en, efter det Udvortes at dømme,
kampdygtig, og forsaavidt til Kamp forpligtet, Individu\-%
alitet kan have sine gode og tilstrækkelige Grunde, hvor\-%
for han netop beslutter sig til ikke at gaae i Kamp.
De personlige Grunde kunne endog være af den Natur,
at han, uagtet han i individuel Selvforstaaelse er sig
deres Gyldighed fuldkommen bevidst, dog føler sig for\-%
pligtet til ikke at meddele dem til Nogen, men resolut
at tage Ansvaret paa sig og finde sig i at blive udsat for
Miskjendelse. Ja, det er ikke engang sagt, at den Mis\-%
kjendte fortier sine Grunde, fordi han bærer paa en eller
anden interessant novellistisk Hemmelighed; den ethiske
Forklaring kan ligefrem være Bevidstheden om, at hans
personlige Grunde enten for stedse eller dog idetmindste
for en Tid vilde blive misforstaaede af Andre. Fra et
andet Synspunkt viser sig noget Lignende. Sæt, at en Kæm\-%
pende er i Kampens Hede stillet paa en høi Post og
betroet et vigtigt Hverv. Opgaven er da at gjøre, ikke
Pligten i theoretisk Almindelighed, thi den kan Ingen
gjøre --- det skulde da være det rene Metaphysikmen\-%
neske --- men \emph{sin} Pligt, altsaa den enkelte bestemte Pligt.
Opgaven er: at lyde sin Ordre og i et afgjørende Øie\-%
blik bruge sin Conduite. Nu kan Situationen, som be\-%
kjendt, let stille sig saaledes, at dersom han virkelig skal
handle paa Conduite, saa kan han ikke lyde sin Ordre, eller,
hvis han punktlig skal lyde sin Ordre, maa han staae
stille og see paa, at Slaget tabes. Saa gjælder det for
Alvor om at fatte en Beslutning: her er Noget at frygte,
og her er Noget at haabe; her skal det vise sig, om
Enkeltvillien --- ikke Metaphysikmenneskets metaphysiske,
men det virkelige Menneskes existerende Enkeltvillie ---
har Charakteer nok til at tage et Ansvar paa sig.

% --- p. 36 ---
\opage{36}\setcounter{footnote}{0}%
Men Charakteer! Charakteer er dog vel ikke et theo\-%
retisk Noget, der i en objectiv Theorie bliver sig selv
objectivt gjennemsigtigt? Nei, personlig er Charakteren
af en saa subjectiv praktisk Natur, at den til Be\-%
slutningen selv alene kan bruge den praktiske Erkjen\-%
delse. Kan en under Indtryk af et stort Ansvar tagen
Beslutning maaskee deduceres theoretisk af Hr. Prof.
Bröchners Theorie: „Jeg erkjender Pligten --- og det
er atter: jer erkjender den theoretisk....; jeg beslutter
% "jer erkjender": so printed (r clear at 600 ppi, both witnesses agree);
% the same quotation reads "jeg erkjender" on p. 29. Transcribed as printed.
mig til at existere i det Erkjendte, og idet Beslutningen
er virkeliggjort, er Indholdet, som realiseret i den en\-%
kelte Villie, væsentlig det samme“? Naar det gjælder
om virkelig Beslutning og virkeligt Ansvar, synes det
endda ikke just at være saa ufornuftigt at gjøre op\-%
mærksom paa Modsætningen mellem theoretisk og prak\-%
tisk Erkjendelse.

Det er nu vistnok ikke Hr. Prof. Bröchners Agt
at ophæve al Forskjel mellem Erkjendelse og Villie;
tvertimod! han mener endog, hvad Modsætningen an\-%
gaaer, at kunne gjøre mig en Indrømmelse. Men hvad
han paa det Bestemteste fastholder, er: at denne Mod\-%
sætning ikke angaaer Indholdet, men alene Formen.
„Begge Former: Erkjendelse og Villie, blive da primi\-%
tive Udtryk for det i Selvbevidstheden som Eenhed
aabenbarede Væsens Energie, for den Selvets Selvbe\-%
stemmen, gjennem hvilken det, bestemmende sig, reali\-%
serer sig: i begge Functioner hævder Selvheden og Ener\-%
gien sig under eiendommelige Grundformer“ (S. 135).
Hvorledes det nærmere forholder sig med disse eien\-%
dommelige Grundformer og de til samme svarende Func\-%
tioner, hvorvidt Enkeltvillien ved disse Functioner vil
være istand til at bringe det til en virkelig Beslutning,
kan sees af Udtalelser som disse: „Den til sit Maal
naaede Erkjendelse vilde, som Udtryk for det i sig
% --- p. 37 ---
\opage{37}\setcounter{footnote}{0}%
gjennemsigtige Væsens Energie, ifølge det ene Væsens
nødvendige Væren i begge Former, umiddelbart omsætte
sig i Villiens, ved Væsenserkjendelsen, Væsenets Selv\-%
aabenbarethed, bestemte Energie, og den Villie, der med
den intensiveste Energie, i den meest omfattende Betyd\-%
ning, var Væsensvillen: Realisation af Væsenet efter dets
% "Væsensvillen": so printed (image and both witnesses), not "-villien";
% within Nielsen's quotation from Brøchner; left as printed
Concretion, vilde være en Villie, der umiddelbart om\-%
satte sig i Bevidsthed, i Erkjendelse, og i en Erkjen\-%
delse, der var bestemt ved Videns almene objectivgyldige
Love.“

Dersom der i disse og lignende Udtalelser skjuler
sig Noget, der virkelig fortjener Anerkjendelse, da maa
det være noget „Fiint og Aandfuldt“, som jeg med min gro\-%
vere Sands ikke er istand til at skimte; thi, saavidt jeg
kan forstaae, indeholdes heri adskillige endog meget ufor\-%
døielige Elementer. For det Første er en psychologisk Ad\-%
skillelse af to for Selvhedsenergien eiendommelige For\-%
mer, der ikke skulle berøre Indholdet, og det uagtet de
udtrykkelig modsættes hinanden som primitive Former,
som „Grundformer“, i lige Grad ulogisk og upsycho\-%
logisk. En Form, der kan forandres, uden at Forandrin\-%
gen væsentlig angaaer Indholdet, er en uvæsentlig, for
Indholdet selv tilfældig Form, og kan umulig være nogen
Grundform. Thi hvad er en Grundform ligefrem efter
sit Begreb? Dog vel en af Væsensindhold opfyldt Form,
en Form, hvis primitive Eiendommelighed ikke kan for\-%
andres uden at Indholdet forandres. Saavidt om det Lo\-%
giske; nu om det Psychologiske. Den Pligt f. Ex.,
at kæmpe for Fædrelandet, kan, idet den erkjendes,
ganske vist opfattes og fastholdes for sig i „Almindelig\-%
hedens og Idealitetens“ Form, og derpaa under Opfyl\-%
delsen omsættes i Villiens, d. e. i „Enkelthedens og Rea\-%
litetens Form“. Men kan denne Omsætning nu siges at
være en Formforandring, der ikke forandrer Indholdet?
% --- p. 38 ---
\opage{38}\setcounter{footnote}{0}%
Man betragte engang den \emph{samme bestemte} Pligt \emph{før} og
\emph{efter} Opfyldelsen. Før Opfyldelsen er Handlingen kun
en Mulighed, dens Bestemmelser Mulighedsbestemmelser
og forsaavidt Almeenbestemmelser. Men Handlingen er
jo Pligtens Indhold; Pligtindholdet er altsaa før Opfyl\-%
delsen et Indhold i Mulighed, et Mulighedsindhold. Men
efter Opfyldelsen, hvor er da Alt forandret! Den Tappre
har opfyldt sin Pligt, han har vovet sit Liv. Idet han
fast besluttet vovede sit Liv, lod han ikke sin existerende
Enkeltvillie som en af de mange Muligheder flyde hen i
Mulighedernes almindelige Strøm; thi den, der vover
saaledes, er ikke en Helt, men en Beruset, der rives med
i Tummelen. At vove er, ethisk seet, at beslutte; at
beslutte er at vove. Idet den Besluttede sætter Pligt\-%
handlingen som det Afgjørende, hæver han sig ikke med
stoisk Ligegyldighed over det mulige Udfald --- den
stoiske Ligegyldighed er en psychologisk Svaghed, der
forvansker det Ethiske --- men han lader Tilskikkelserne
komme over sig, at Pligtopfyldelsens \emph{mulige} Indhold
kan ved Handlingen, den enkelte virkelige Handling,
blive et \emph{virkeligt} Indhold. I den virkelige Handling
har Villien vovet sig selv som existerende Enkeltvillie,
har ved at sætte \emph{sin Existens} ind i Indholdet væsentlig
forandret Indholdet. At Hr. Prof. Bröchner har kunnet
oversee en saa iøinefaldende Indholdsforandring, er be\-%
synderligt.

Endnu mere besynderlig er den Anstrengelse, Hr.
Professoren gjør, for at forklare „det i Selvbevidstheden
som Eenhed aabenbarede Væsens Energie“, „den Selvets
Selvbestemmen, gjennem kvilket det, bestemmende sig,
% "kvilket": so printed (k at 600 ppi, both witnesses agree), for "hvilket";
% the same quotation reads "gjennem hvilken det" on p. 36. Printer's error,
% transcribed as printed.
realiserer sig“ i de to Grundformer: „Erkjendelse og
Villie“. Forklaringen er virkelig curieus. For at hævde
Erkjendelsens --- den objective Videns --- „Princip“
overfor Villien og sørge for, at „det subjective Væsen“
% --- p. 39 ---
\opage{39}\setcounter{footnote}{0}%
kan komme til at virke „\emph{efter sin Heelhed}, sin udeelte
Personlighed“, etablerer den skarpsindige Psycholog ikke
mindre end tre principielt selvstændige Selvhedsenergier
i det ene og samme personlige Væsen. Den første af
disse Energier kjendes væsentlig derpaa, at den i sig
har Indhold, men ingen Form, de to andre udmærke sig
ved, at de hver i sig ere Former, eiendommelige „Grund\-%
former“ uden Indhold. At „det i Selvbevidstheden som
Eenhed aabenbarede Væsens Energie“ i sig selv er form\-%
løst, sees udtrykkelig deraf, at det, for at erholde Form,
maa vende sig først til Erkjendelsen, derpaa til Villien;
at Erkjendelse og Villie, hver for sig, maae opfattes som
rene Former uden Indhold, sees udtrykkelig deraf, at
disse Former ikke paa nogen Maade ere istand til at
forandre Indholdet.

Tør vi nu tage denne skarpsindige Inddeling væsent\-%
lig, efter de angivne Begrebsbestemmelser bogstave\-%
lig, saa komme vi da rigtignok til en Dualisme, der
er ganske anderledes haard end den af mig paaviste.
Er virkelig Erkjendelsen i den Forstand uafhængig af
Villien, at den faaer sin Selvhedsenergie, ikke fra Vil\-%
lien som Villie, men fra den første, af Villien uafhæn\-%
gige Selvbestemmen, hiin endnu formløse, mystiske Potens,
der oprindelig søger sin Form i Erkjendelsen: saa have
vi jo en heel Vidensregion, en Kreds af objective Theo\-%
rier, om hvilken det med Sandhed kan siges, at den
staaer uafhængig af Villien, ganske uafhængig saavel
med Hensyn til Indhold som til Form. Thi den objec\-%
tive Viden kan da umulig afhænge af Villien efter sit
Indhold, naar Indholdet kommer fra den formløse Potens;
heller ikke kan Viden afhænge af Villien efter sin Form,
naar Erkjendelsesformen ikke blot har en ligesaa stor,
men en endnu større Primitivitet end Villiesformen.
% --- p. 40 ---
\opage{40}\setcounter{footnote}{0}%
Dersom dette ikke er en psychologisk Dualisme, hvad
er saa Dualisme?

Maaskee tør vi ikke lægge en saadan Vægt paa
Adskillelsen af Form og Indhold. Vi ville da forsøge
det paa en anden Maade og antage, at den hele Disting\-%
veren mellem tre forskjelligartede Selvhedsenergier ikke
har stort at betyde; thi ved at antage dette synes vi,
idetmindste paa nogle Steder, at komme Forfatterens
Mening nærmere. Det siges os, at „den til sit Maal
naaede Erkjendelse“, vilde umiddelbart omsætte sig i Vil\-%
liens, ved Væsenserkjendelsen bestemte Energie; det siges
os, at „den Villen, der med den intensiveste Energie var
Væsensvillen“, vilde være „den Villie, der umiddelbart
omsatte sig i Bevidsthed, i Erkjendelse“, og det oveni\-%
kjøbet i en Erkjendelse, der var bestemt ved Videns al\-%
mene, objectivgyldige Love. Af hvad her saaledes siges,
faae vi en Forestilling om, at det er, „Mediationen“, den
gamle velbekjendte speculative Mediation, der skal op
igjen og holde Tridt med den raisonnerende Kritik.
Men Kritiken og Mediationen kunne ikke forliges. Skal
Mediationen raade, gaaer det paa ingen Maade an at
antage en væsentlig Formforandring uden tilsvarende Ind\-%
holdsforandring. Man sige om Mediationen, hvad man
vil: i Kapitlet om Form og Indhold tillader den ingen
Halvering; dens Tænkning er immanent, og den immanente
Tænkning indlader sig ikke paa kritisk Udstykning af
Form og Indhold. Hvorledes skulle vi vel kunne tænke
os, at den til sit Maal naaede Erkjendelse vilde bære
sig ad med at omsætte sig --- umiddelbart --- i Villie?
Erkjendelsens Maal er jo lutter objectiv Theorie, lutter
Videnskab. Saalænge den objective Erkjendelse ikke
har udtømt alt muligt Vidensstof, maa den naturligviis
gaae frem fra Viden til Viden uden at omsætte sig i
Villie. Da nu Forfatteren ikke vil indrømme nogen
% --- p. 41 ---
\opage{41}\setcounter{footnote}{0}%
absolut Vidensgrændse, vil Erkjendelsen selvfølgelig al\-%
drig naae sit Maal, og den forventede Omsætning aldrig
indtræde. Sæt, at Erkjendelsen, efter at have naaet et
relativt Maal, omsætter sig til Villie: hvad standser da
Erkjendelsen; hvad er det, der bevirker Omsætningen?
Ifølge Mediationen skulde Erkjendelsen standse ved sig
selv, ved sig selv gaae over til Villie, fordi den i sit
Væsen er Villie. Men det forbyder Kritiken; den rai\-%
sonnerende Kritik har udtrykkelig sat Erkjendelsesformen
for sig, uafhængig af Villiesformen. Ifølge sund Menne\-%
skeforstand skulde det være Villien, der standsede Er\-%
kjendelsen og bevirkede Omsætningen; men det strider
imod den opstillede Theorie, der sætter Erkjendelsen uaf\-%
hængig af Villien. Det maa altsaa være hiin i sig form\-%
løse Selvhedsenergie, hiin mystiske Potens, der hverken
er Erkjendelse eller Villie, som maa omsætte Erkjendelsen
i Villie. Men selv naar dette antages, er her endnu
ingen Grund til, at Erkjendelsen, „som Udtryk for det i
sig gjennemsigtige Væsens Energie, ifølge det ene Væsens
nødvendige Væren i begge Former“, skulde „umiddel\-%
bart omsætte sig“ i Villiesformen. Den eneste Grund,
Forklaringen indeholder, er en Slags Grund til at „det
i sig gjennemsigtige Væsens Energie“, efter at have ud\-%
trykt sit Indhold i Erkjendelsens Form, kan bytte om
og udtrykke sit Indhold i Villiens Form. Men et saa\-%
dant Formskifte er da saare langt fra at betegne den
ene Forms Overgang i den anden.

Forsøge vi at tænke Overgangen ved at gaae ud
fra Villiens Selvomsætten i Viden, saa ere vi lige nær.
Hvad er det for en Villie, der --- med „den intensiveste
Energie“ --- umiddelbart omsætter sig i Bevidsthed?
Ifølge Forfatterens egen Antagelse kan Villien jo som
Villie ikke være uden Bevidsthed; Villien med den „in\-%
tensiveste Energie“ maa altsaa være en Villie med den
% --- p. 42 ---
\opage{42}\setcounter{footnote}{0}%
intensiveste Bevidsthed. Men naar Villien med den in\-%
tensiveste Energie selv, nemlig i Villiens Grundform,
\emph{er} den intensiveste Bevidsthed, hvorledes kan den da
\emph{omsætte} sig i Bevidsthed, i det, den allerede er?
Her er jo ingen Omsætning. Skal den vidunderlige Om\-%
sætning maaskee bestaae deri, at Villien, efter at den
med en mindre intensiv Energie en Tidlang har holdt
sig i Form af ordinær Villie, simpel, praktisk Charak\-%
teervillie, omsider naaer den intensiveste Energie, og saa
pludselig ved Hjælp af den intensiveste Energie springer
om, --- „umiddelbart“, som naar det hedder, at Vinden
springer om --- og forvandler sig fra Villie til Ikke-%
Villie, til Erkjendelse, til objectiv Erkjendelse, Erkjen\-%
delse af Videns almene objectivgyldige Love: da vil den
hele Forklaring nærme sig stærkt til det Meningsløse.
Den Philosoph, der i et bestemt psychologisk Exempel
skulde kunne tydeliggjøre os en saadan Omsætning, en
saadan umiddelbar objectiv Springen om i det Subjective
maatte være en sand theoretisk Hexemester, en objectiv
psychologisk Tusindkunstner.

En sidste Udvei staaer endnu tilbage; heller ikke
den tør lades uforsøgt. For at hævde Personlighedens
udeelte Eenhed med sig selv i Erkjendelse og Villie,
beraaber Hr. Professoren sig paa hvad han kalder „det
i Selvbevidstheden som Eenhed aabenbarede Væsens
Energie“, paa „den Selvets Selvbestemmen, gjennem
hvilken det, bestemmende sig, realiserer sig“ o. s. v.
Det er altsaa Forholdet imellem Villien og Selvets Selv\-%
bestemmen, hvorpaa Opmærksomheden maa fæste sig.
Ifølge den sædvanlige Synsmaade kunde man fristes til
at antage, at Selvets Selvbestemmen egentlig var dets
oprindelige Villen, og Villien selv Personlighedens Cen\-%
trum; men dette er ingenlunde Tilfældet. „Det i Selv\-%
bevidstheden som Eenhed aabenbarede Væsens Energie“
% --- p. 43 ---
\opage{43}\setcounter{footnote}{0}%
har, naar vi abstrahere fra Erkjendelsesenergiens pseu\-%
doselvstændige Selvbestemmen --- mirabile dictu --- to
Slags særlig bestemmende Selvbestemmen, „gjennem
hvilken det, bestemmende sig, realiserer sig“. Den første
Selvbestemmen er den i sig formløse --- den mystiske
Potens; den sidste, den i sig indholdsløse, er Villien.
Villien som Villie er altsaa ikke Personlighedens selvbe\-%
stemmende Eenhed; thi Villien fremkommer først, naar
den første, formløse Selvbestemmen bestemmer sig til at
lægge Bevidsthedsindholdet ind i den anden bestemmende
Selvbestemmen, den nemlig, der bestemmer Formen.
Villien, seet i denne Belysning, er da saa godt som util\-%
regnelig. Villien, ak den stakkels Villie, er bleven fangen
af en mystisk Potens og kan ikke oprindelig bestemme
sig selv: dens Selvbestemmelse er hver Gang afhængig
af, hvorvidt den Selvets første Selvbestemmen, „gjennem
hvilken det, bestemmende sig, realiserer sig“, vil be\-%
stemme den til Selvbestemmen.

Det er paa et saadant Grundlag, at Hr. Professor
Bröchner vil hævde Personlighedens udeelte Eenhed!
Først gjør han af egen theoretisk Magtfuldkommenhed
den existerende Bevidsthed til en videnskabelig Theorie;
dernæst confunderer han det Almenes subjective Gyldig\-%
hed med den objective, endog i den Grad, at han ikke
seer sig istand til at gjøre Forskjel imellem Videnskabs\-%
grunde og Samvittighedsgrunde, imellem en Erkjendelse
af Pligtloven og en Erkjendelse af Tyngdeloven. Der\-%
næst, efter at det i Kjød og Blod existerende Aands\-%
menneske er blevet forvexlet med et Skyggerids af et
ved Abstractioner afbleget Metaphysikmenneske, bliver
Enkeltvilliens Almeenbegreb af en Feiltagelse sat i den
for sin Selvstændighed kæmpende Villies Sted. Endelig,
efter at det subjective Universalvæsen med de tre Selv\-%
hedsenergier, den ene med Indhold uden Form, de to
% --- p. 44 ---
\opage{44}\setcounter{footnote}{0}%
andre med Form uden Indhold, er blevet sig selv saa
objectivt gjennemsigtigt, at Subjectivitetens oprindelige
Selvbestemmen kan i fuldkommen Gjennemsigtighed be\-%
stemme sig til Villiens Selvbestemmen, er Maalet naaet:
Handlingen er Theorie og Theorien Handling, Forskjellen
i al Fald kun en Formforskjel, der ikke forandrer Indholdet.
Paa et saadant Grundlag er det nu en let Sag at negte
Videns absolute Grændse og forkaste Mysteriet; paa et
saadant Grundlag kan en kritiserende Philosoph da
sagtens bevise baade Umuligheden af Guds Personlighed
og Umuligheden af Bønnens Gyldighed og Umuligheden
af Sjælens Udødelighed.

Paa et saadant Grundlag, et luftigt Underlag af
væsenstomme Abstractioner, skjæve Deductioner, ud\-%
tværede Vendinger og stakaandede Bemærkninger har
Hr. Professor Bröchner nu virkelig styltet sig op til en
philosophisk Myndighed, en Videns Ypperstepræst, der
renser Troen, bevarer „det Religiøse“ ved at „afskaffe
Religionerne“ og forkynder Metaphysikmenneskets, det
„universelle Menneskevæsens“ Evangelium, idet Viden
vælges. Saa lyde hans Ord: „Idet Viden vælges, vælges
nemlig det, der, som Væsensudtryk for det universelle
Menneskevæsen, i sig bærer Muligheden til at afslutte
sig til en Totalitet med indre Uendelighed; det der, idet
det har det ideelle Principat i den udelelige menneske\-%
lige Aand, kan blive Normen for en ethisk Livsheelhed;
det, der, medens det gjennemskuer de endelige Religions\-%
formers Væsen, begriber dem i deres Udspring af den
sig udfoldende, hen mod sit Princip stræbende Menne\-%
skenatur, og derfor ikke i dem beholder en forstyrrende
Bevidsthedsskranke, og som, idet det i de historiske Re\-%
ligionsformer udskiller det Forgængelige og Endelige fra
det Blivende, Væsentlige og Evige, det Ikke-Religiøse
fra det Religiøse, kan staae i Harmonie med og i sig
% --- p. 45 ---
\opage{45}\setcounter{footnote}{0}%
optage dette Væsentlige, det Religiøse efter dets sande og
evige Betydning.“ (S. 225---26).

Dersom Nogen nu, efter møisommeligt at have ar\-%
beidet sig igjennem Hr. Professorens Skrift: \emph{Problemet
om Tro og Viden}, erklærede reent ud, at Grunden,
hvorfor Hr. Prof. Bröchner ved denne Leilighed er op\-%
traadt mod det sidste mulige Forsøg paa at værge Aaben\-%
baringstroen mod Videnskabens Angreb, aabenbart ikke
er Kjærlighed til Videnskaben, men et ulmende Had til
den positive Christendom, mon man da ikke vilde tage
kraftig til Gjenmæle og opløfte et Ramaskrig: „Ingen
Insinuationer, ingen Personligheder! prøver Beviserne,
men skumler ikke over Bevæggrundene!“ Jeg antager,
at man vilde gjøre det, og med Rette. Men hvad be\-%
viser det? Det beviser jo netop, hvad jeg gjør gjæl\-%
dende, men hvad Hr. Prof. Bröchner ikke vil indrømme,
at der er en Forskjel, en afgjort principiel Forskjel mel\-%
lem Vidensspørgsmaal og Samvittighedsspørgsmaal, mel\-%
lem Videnskabsgrunde og Personlighedsgrunde. Det
beviser just, hvad jeg gjør gjældende, men hvad Hr.
Prof. Bröchner udtrykkelig benegter, at der er en For\-%
skjel, en uimodsigelig Principforskjel imellen den klare % sic: "imellen" (for imellem) as printed
Indsigt i en objectiv Theorie og „Speilingen“ af det
Personlighedens Indre, som aldrig gaaer op i noget
videnskabeligt Begreb.

To Vidende kunne blive hinanden objectivt gjennem\-%
sigtige i den samme Viden af den samme Videnskab;
men to Personligheder kunne aldrig blive hinanden objec\-%
tivt gjennemsigtige i de samme personlige Anliggender;
thi dette staaer fast: jo dybere Personlighed, desto rigere
Inderlighed; men ingen Videnskab eier den magiske Nøgle,
der lukker Inderlighedens Dør op.

Men endskjøndt der i de existerende Personligheder
er Dybder, hvortil Videns objective Lysning aldrig naaer
% --- p. 46 ---
\opage{46}\setcounter{footnote}{0}%
ind, saa følger dog ingenlunde deraf, at der ikke skulde
kunne fældes en almindelig, tilnærmelsesviis begrundet
Dom om personlige Handlinger, om virkelige Charak\-%
terer og deres større eller mindre Overeensstemmelse saa\-%
vel med sig selv som med det Almeenmenneskelige;
kun at Opfattelsen, som sagt, bliver en phychologisk % sic: "phychologisk" as printed
Speiling ikke en udtømmende Begriben. Med Hen\-%
syn hertil kan jeg da meget vel, uden at gaae det Per\-%
sonlige for nær, gjøre opmærksom paa, at den Interesse,
der i Hr. Prof. Bröchners foreliggende Kritik, træder i
Forgrunden, ikke er en positiv Vidensinteresse, men en
negativ Troesinteresse. Jeg skjønner det af følgende
Indicier. Dersom det virkelig havde været Hr. Prof.
Bröchners Hensigt at gjøre Videnskabens Ret gjældende
imod Overgreb fra Troens Standpunkt: da maatte han, i
dette Punkt idetmindste, have sluttet sig til mig. For\-%
handlingen imellem os var da bleven en videnskabelig
Forhandling om noget reent Videnskabeligt, f. Ex. om
Problemerne i Grundideernes Logik. Af dette Arbeide
foreligger en Række Analyser paa ikke mindre end 60
Ark. Her maatte altsaa, hvorledes Resultatet end blev,
være Noget at undersøge. At jeg ved Siden af disse
reent videnskabelige Bestræbelser desuden mente at kunne
bevare og forsvare den gamle Aabenbaringstro, kunde Hr.
Prof. jo indtil videre have ladet staae hen som et i hans Øine
psychologisk Curiosum, en „Speiling“ af noget Person\-%
ligt, der ikke ret lod sig udtømme i det objective Be\-%
greb, idetmindste ikke, førend jeg var naaet saavidt, at
jeg kunde fremlægge min Religionsphilosophie som et
gjennemarbeidet Hele. Men nu er Hr. Prof. Brøchner
netop gaaet den omvendte Vei. Langt fra at vente paa
en authentisk og sammenhængende Fremstilling af min
Troeslære, har han med stort Hastværk kastet sig over
„Bemærkninger“, „Referater“ og „aphoristiske Udta\-%
% --- p. 47 ---
\opage{47}\setcounter{footnote}{0}%
lelser“, har behandlet hver Sætning, hvert Ord, hver
Tøddel, urgerende, præciserende, distingverende, re\-%
ducerende, concederende, restringerende, raisonerende
o. s. v. med en Iver saa systematisk, som havde
han et fuldfærdigt System for sig. Dette Hastværk,
denne Iver tyder aabenbart hen paa, at Hovedopgaven
er at komme „Aabenbaringstroen“ tillivs og faae „Reli\-%
gionerne“ afskaffede for at kunne bevare det „Religiøse“.

At Vidensinteressen ikke staaer i Forgrunden, kan
jeg let bevise. Jeg behøver jo blot at henvise Læseren
til den Bedømmelse af Grundideernes Logik, der findes i
Hr. Professorens Skrift (S. 154---160). Til en begrundet
Dom om denne Bedømmelse er det aldeles ikke for\-%
nødent, at den dannede Læser skal være synderlig inde
i Metaphysiken og vide Besked enten om Forholdet
imellem System og Schema eller om „Magten“ eller om
„det Antilogiske“; naar han blot har en nogenledes klar % a hairline mark stands over the a; read as a speck, like the one over "en" on p. 28 (collation 2026-09-22; burden of proof on the accented sort; JBIG2, no sort claim)
Forestilling om, hvad en Forfatter, der idetmindste har
lagt an paa at præstere noget Videnskabeligt, er beret\-%
tiget til at fordre af den, som vil underkaste hans Ar\-%
beide en videnskabelig Prøvelse, vil han uden Vanske\-%
lighed fatte, hvad Enhver, der har Agtelse for Videnskab,
maa dømme om Hr. Prof. Bröchners videnskabelige Be\-%
dømmelse. Thi det er vist, at et halvmetaphysisk Fruen\-%
timmer, naar hun blot forstod at bruge saadanne Ord
som „Subreption“, „Tilsnigelse“, „Uklarhed“, „Vilkaar\-% image: the opening „ shows as a comma followed by a dot
lighed“, „Sophistik“, „Metaphor“, „Kategorie“, „ikke de\-%
duceret“ o. dsl., paa Steder, hvor de klinge, vilde i
en Avisartikel kunne levere en ligesaa grundig, ligesaa
videnskabelig Bedømmelse, og saa ovenikjøbet gjøre den
nok saa læselig.

At Grundideernes Logik ikke falder i Hr. Professor
Bröchners Smag, maa vistnok forklares af flere Grunde,
blandt Andet deraf, at den slet ikke ligger for hans kri\-%
% --- p. 48 ---
\opage{48}\setcounter{footnote}{0}%
tiske Talent. „Det er en gammel Indvending, at Philo\-%
sophien hænger sig i Ord, at dens Forklaringer ere Ord\-%
forklaringer, dens Stridigheder Ordstrid, dens Viisdom
Ordkram. Denne Indvending er ikke ganske uden Grund,
naar man eensidig fæster Opmærksomheden paa Skygge\-%
siderne, paa de ufrugtbare Subtiliteter, de smagløse Ud\-%
skeielser i tom Ordkløgt, om hvilke en Johan af Salis\-%
bury, en Erasmus fra Rotterdam og hos os en Holberg
have, hver i sin Tidsalder, holdt Rettergang med den
speculative Videns Clienter.“ Med disse Ord begynder
anden Deel af den omhandlede Logik, og maaskee er der i
dem noget Frastødende. Vistnok lægges der Vægt paa Ordet.
„Philosophien er Ordets Videnskab; den maa nødvendig
holde paa Ordet, saasandt den vil holde paa Tanken og
fastholde Begrebet. Ligesom en Sjæl ikke kan være
uden et Legeme, saaledes kan Begrebet som det i sig
bestemte Begreb heller ikke være uden et bestemmende
Ord; thi dersom Tanken ikke var, saa var der intet
Ord, og dersom Ordet ikke var, saa var der heller ingen
Tanke.“ Men under den videre Udvikling bliver det
dog med afgjørende Eftertryk fremhævet, at den af
Ordet betingede Tankeudvikling altid bør gaae i Eet
med en tilsvarende Sagudvikling. Her er nu, hvis jeg
ikke feiler, det egentlige Collisionspunkt mellem Grund\-%
ideernes Logik og Hr. Prof. Bröchners Talent.

At Hr. Prof. Bröchners Talent er et kritisk Talent
af en egen Art, nemlig et Talent til philosophisk Ord\-%
kritik, skal i denne Sammenhæng ikke henkastes som en
flygtig Paastand, men til en Afslutning anskueliggjøres
ved et Par slaaende Exempler. For at belyse, hvad han
kalder „Theoriens Conseqventser med Hensyn til det
Psychologiske“ (S. 178 flgd.), begynder Hr. Professoren
med et Par udførlige, fra mit Skrift: „Om den gode Villie“
hentede Citater, og indleder saa Belysningen selv med
% --- p. 49 ---
\opage{49}\setcounter{footnote}{0}%
følgende Erklæring: „Vi have gjengivet disse Udviklinger,
ogsaa de mod Slutningen anbragte „lazzi“, saa fuldstæn\-%
dig, fordi kun en saadan Fuldstændighed kan give et
tydeligt Billede af den Maade, paa hvilken Theorien
maa forsvares“ (S. 183). Ved dette forcerede Tilløb skal
Læseren naturligviis imponeres til at antage, at hvad jeg
har sagt ikke blot er noget flaut og intetsigende Tøieri, men
ovenikjøbet noget høist Upassende. Mine Ord ere disse:
„Staaer det da fast, at Troen paa ingen Maade kan for\-%
vandles til Viden eller Viden til Tro, saa staaer det
ogsaa fast, at Troen ligesaa umuligt kan halveres ved
Viden som Viden ved Tro. Kan en Bevidsthed ikke
halveres af Venskab og Mathematik, saa kan den
endnu mindre halveres af Ethik og Astronomie. Der
gives mathematiske Venner, men der gives ingen ven\-%
skabelig Mathematik; der gives forelskede Chemikere,
men ingen chemisk Forelskelse, ingen forelsket Chemie.
Der kan gives ethiske og uethiske Astronomer, men
Astronomien er hverken ethisk eller uethisk. Viden\-%
skaben er med andre Ord hverken ethisk eller uethisk,
hverken troende eller vantro; Kjærlighed er praktisk,
men Videnskab er theoretisk.“ Det er kun faa Linier,
hvorom her handles, og dog indrømmer jeg, at den, der
kan rokke den i disse faa Linier sammentrængte Grund\-%
tanke, naar det gjælder en Sagudvikling, han kuldkaster
hele min Anskuelse. „Men“ --- indvender man maaskee
--- ;„Hr. Prof. Brøchner har jo rokket den!“ Jeg svarer: % image: a ;-shaped mark before „, transcribed as printed (stray sort or speck; JBIG2, no sort claim)
Hr. Professor Bröchner gaaer, hvad enhver nogenledes
opmærksom Læser kan see, aldeles ikke ind paa Sagen;
nei, han lirker ved Ordene: hans philosophiske Kritik er
i dette som i andre meget vigtige Stykker en, endog tem\-%
melig pusillanim, Ordkritik. Man læser (S. 187): „De
„livlige“ Analogier, N. tilføier til Oplysning, ere nærmest
kun et lidet heldigt Forsøg i det Burleske, og bevise,
% --- p. 50 ---
\opage{50}\setcounter{footnote}{0}%
med Hensyn til den omhandlede Sag, om muligt, mindre
end Intet“.\footnote{Havde Hr. Prof. Bröchner med Hensyn til „de mathematiske
Venner“ og „den venskabelige Mathematik“ virkelig oplyst eller
endog blot gjort Forsøg paa at oplyse Noget, der angik „den
omhandlede Sag“, da havde vi maaskee faaet en interessant Op\-%
lysning om, hvorledes det personlige Forhold mellem mathematiske
Venner kan bidrage til \emph{enten at svække eller forstærke
de objective, mathemathiske Beviser}, hvorledes, med % sic: "mathemathiske" as printed
andre Ord, Personlighedsgrunde kunne forvandles til Viden\-%
skabsgrunde, en Oplysning, der sikkerlig ikke vilde undlade
at bevirke en stor Revolution saavel i Mathematiken som
i de øvrige Videnskaber. Men istedenfor at oplyse Noget
om Sagen, har Hr. Professoren, maaskee af personlige
Grunde, foretrukket at oplyse Noget om Ordene. Oplysningen
er, „at man ikke uden videre, fordi et Adjectiv kan prædiceres
om et Substantiv, derfor omvendt, kan prædicere et til Substan\-%
tivet svarende Adjectiv om et til Adjectivet svarende Substantiv“,
en Oplysning, som, hvad Sagen angaaer, omtrent staaer paa lige
Trin med Per Degns Oplysning om Imprimatur: „Alle de Ord,
som nævnes kand, henføres til 8 Ting, som er Nomen, Prono\-%
men, Verbum, Participium, Conjugatio, Declinatio, Interjectio.“

I hvilken Grad Hr. Professor Bröchner kan forglemme
Sagen, blot for at kritisere Ordene, og hvor streng han er,
naar han kritiserer Ordene, sees deraf, at hans Skarpsindighed
endog i de anførte Exempler har kunnet opdage „det Burleske“.
Ærbarhed er en Dyd, selv i en philosophisk Fremstilling. Men
da Grændsen for det Ærbare jo ikke er objectiv videnskabelig,
kunne de philosopherende Personligheder dog ogsaa drive Ær\-%
barhedsfordringen saavidt, at det gaaer ud over Fremstillingen.
Seer Hr. Professoren sig istand til i „en theoretisk udtømmende
Begrebsbestemmelse“ at bevise, at jeg i de anførte Exempler har
indblandet „urene Bestanddele“, der kunne henføres under Kate\-%
gorien af „det Burleske“, da maa jeg naturligviis bøie mig for
objective, uimodsigelige Grunde; men indtil et saadant Beviis
fremlægges, forbeholder jeg mig min personlige Ret. Thi det
er virkelig min Overbeviisning, at Philosophien, selv om den
% note continues on p. 51
var personificeret i det rene Metaphysikmenneske og ved denne
Forvandling --- der naturligviis ikke maatte vedrøre Indholdet, men
blot Formen --- bleven saa ærbar som en gammel Skolejomfru:
saa skulde den gamle Jomfru, medmindre hun fik et hysterisk
Anfald af Arrighed, dog i mine uskyldige Exempler have ondt
ved at finde „det Burleske“.}

% --- p. 51 ---
\opage{51}\setcounter{footnote}{0}%
Selv det lille Stykke (S. 184---186), hvor den kriti\-%
serende Forfatter omsider efter adskillige Omsvøbsbe\-%
mærkninger lader, som om han gik ind paa Sagen, ved\-%
bliver Kritiken ligefuldt at være en Ordkritik, hvori den
ene skjæve Bemærkning afløser den anden. „Idet Viden
anerkjender sin Grændse, anerkjender den i det, der er
Troens Sphære, et Mysterium for Viden. Idet Troen
har en Grændse mod Viden, bliver det, der er Videns
Gjenstand, et Mysterium for Troen. Troen er altsaa
Videns Mysterium, Viden Troens, og Mysteriet, der,
saaledes opfattet, ikke har noget Mysteriøst, forhindrer
paa denne Maade ikke Indskrænkningen, men udtrykker
den netop.“ Dette Exempel paa tomt Ordspil kan her
figurere som instar omnium; kun gjælder det om at
fatte „det Fine“. „Det Fine“ er, at et Mysterium, naar
man veed, hvad det er: „en Grændse mod Viden“, ikke
længer er et Mysterium; thi man veed jo, hvad det er: „en
Grændse mod Viden“. Altsaa er en Hemmelighed, naar
jeg veed, hvad den er, nemlig det, jeg ikke kan vide,
ikke længer nogen Hemmelighed, thi jeg veed jo, hvad
en Hemmelighed er, nemlig: det, jeg ikke kan vide.

Den følgende Bemærkning er vel ikke slet saa tom,
men den er baade vind og skjæv. „Viden skal, ved
at underforstaae de i \emph{Mysteriet skjulte Mellembe\-%
stemmelser}, kunne tænke sig den Magtaabenbarelse,
„i hvilken den Troende seer Miraklet“, som bestaaende
med Lovenes Uforanderlighed, men dette kan den kun,
% --- p. 52 ---
\opage{52}\setcounter{footnote}{0}%
naar den, medens Miraklets Navn beholdes, lader Miraklet
som \emph{Mirakel forsvinde i et supponeret Ratio\-%
nelt}.“ De Læsere, som ønske at kjende min Opfattelse
af Miraklet i en tilsvarende Sagudvikling, henviser jeg
til mine Forelæsninger om Hindringer og Betingelser
for det aandelige Liv i Nutiden (IX og X); om Be\-%
mærkningen i den foreliggende Ordkritik kan jeg fatte
mig kort. Det for Miraklet Charakteristiske er slet ikke
det, at der gjøres Brud paa Lovene; men det, at en
Kjendsgjerning bliver umiddelbart henført til Guds al\-%
mægtige Villie. At Grundforholdet imellem Naturriget
og Guds almægtige Villie er et indre Uendelighedsfor\-%
hold, altsaa et Forhold, hvori Villien virker i Eet med
Lovene, kan indsees i al Almindelighed, uden at Grund\-%
forholdet selv i sine virkelige Articulationer derfor ophører
at være et saavel for Viden som for Troen absolut Myste\-%
rium. Hvorvidt et saadant Mysterium skal antages eller
benegtes, er ikke et Videnskabsspørgsmaal; nei, det er
endnu som før et Samvittighedsspørgsmaal.

En ukritisk Sammenblanding af Samvittighedsspørgs\-%
maal med Videnskabsspørgsmaal er vel det stærkeste
Vidnesbyrd imod Videnskabeligheden i Hr. Prof. Bröchners
philosophiske Kritik.

\begin{center}\rule{3cm}{0.4pt}\end{center}

\end{document}
